%----------
%   IMPORTANTE
%----------

% Si nunca has utilizado LaTeX es conveniente que aprendas una serie de conceptos básicos antes de utilizar esta plantilla. Te aconsejamos que leas previamente algún tutorial (puedes encontar muchos en Internet).

% Esta plantilla está basada en las recomendaciones de la guía "Trabajo fin de Grado: Escribir el TFG", que encontrarás en http://uc3m.libguides.com/TFG/escribir
% contiene recomendaciones de la Biblioteca basadas principalmente en estilos APA e IEEE, pero debes seguir siempre las orientaciones de tu Tutor de TFG y la normativa de TFG para tu titulación.

% Encontrarás un ejemplo de TFG realizado con esta misma plantilla en la carpeta "_ejemplo_TFG_2019". Consúltalo porque contiene ejemplos útiles para incorporar tablas, figuras, listados de código, bibliografía, etc.


%----------
%	CONFIGURACIÓN DEL DOCUMENTO
%----------

% Definimos las características del documento y añadimos una serie de paquetes (\usepackage{package}) que agregan funcionalidades a LaTeX.

\documentclass[12pt]{report} %fuente a 12pt

% MÁRGENES: 2,5 cm sup. e inf.; 3 cm izdo. y dcho.
\usepackage[
a4paper,
vmargin=2.5cm,
hmargin=3cm
]{geometry}

% INTERLINEADO: Estrecho (6 ptos./interlineado 1,15) o Moderado (6 ptos./interlineado 1,5)
\renewcommand{\baselinestretch}{1.15}
\parskip=6pt

% DEFINICIÓN DE COLORES para portada y listados de código
\usepackage[table]{xcolor}
\definecolor{azulUC3M}{RGB}{0,0,102}
\definecolor{gray97}{gray}{.97}
\definecolor{gray75}{gray}{.75}
\definecolor{gray45}{gray}{.45}

% Soporte para GENERAR PDF/A --es importante de cara a su inclusión en e-Archivo porque es el formato óptimo de preservación y a la generación de metadatos, tal y como se describe en http://uc3m.libguides.com/ld.php?content_id=31389625. En la carpeta incluímos el archivo plantilla_tfg_2017.xmpdata en el que puedes incluir los metadatos que se incorporarán al archivo PDF cuando lo compiles. Ese archivo debe llamarse igual que tu archivo .tex. Puedes ver un ejemplo en esta misma carpeta.
\usepackage[a-1b]{pdfx}

% ENLACES
\usepackage{hyperref}
\hypersetup{colorlinks=true,
	linkcolor=black, % enlaces a partes del documento (p.e. índice) en color negro
	urlcolor=blue} % enlaces a recursos fuera del documento en azul

% EXPRESIONES MATEMATICAS
\usepackage{amsmath,amssymb,amsfonts,amsthm}

\usepackage{txfonts} 
\usepackage[T1]{fontenc}
\usepackage[utf8]{inputenc}

%\usepackage[english]{babel} 
%\usepackage[babel, english=american]{csquotes}
\usepackage[spanish,es-tabla,es-noshorthands]{babel} 
\usepackage[autostyle=true]{csquotes}

\AtBeginEnvironment{quote}{\small}
% diseño de PIE DE PÁGINA
\usepackage{fancyhdr}
\pagestyle{fancy}
\fancyhf{}
\renewcommand{\headrulewidth}{0pt}
\rfoot{\thepage}
\fancypagestyle{plain}{\pagestyle{fancy}}

% DISEÑO DE LOS TÍTULOS de las partes del trabajo (capítulos y epígrafes o subcapítulos)
\usepackage{titlesec}
\usepackage{titletoc}
\titleformat{\chapter}[block]
{\large\bfseries\filcenter}
{\thechapter.}
{5pt}
{\MakeUppercase}
{}
\titlespacing{\chapter}{0pt}{0pt}{*3}
\titlecontents{chapter}
[0pt]                                               
{}
{\contentsmargin{0pt}\thecontentslabel.\enspace\uppercase}
{\contentsmargin{0pt}\uppercase}                        
{\titlerule*[.7pc]{.}\contentspage}                 

\titleformat{\section}
{\bfseries}
{\thesection.}
{5pt}
{}
\titlecontents{section}
[5pt]                                               
{}
{\contentsmargin{0pt}\thecontentslabel.\enspace}
{\contentsmargin{0pt}}
{\titlerule*[.7pc]{.}\contentspage}

\titleformat{\subsection}
{\normalsize\bfseries}
{\thesubsection.}
{5pt}
{}
\titlecontents{subsection}
[10pt]                                               
{}
{\contentsmargin{0pt}                          
	\thecontentslabel.\enspace}
{\contentsmargin{0pt}}                        
{\titlerule*[.7pc]{.}\contentspage}  


% DISEÑO DE TABLAS. Puedes elegir entre el estilo para ingeniería o para ciencias sociales y humanidades. Por defecto, está activado el estilo de ingeniería. Si deseas utilizar el otro, comenta las líneas del diseño de ingeniería y descomenta las del diseño de ciencias sociales y humanidades
\usepackage{multirow} %permite combinar celdas 
\usepackage{caption} %para personalizar el título de tablas y figuras
\usepackage{floatrow} %utilizamos este paquete y sus macros \ttabbox y \ffigbox para alinear los nombres de tablas y figuras de acuerdo con el estilo definido. Para su uso ver archivo de ejemplo 
\usepackage{array} % con este paquete podemos definir en la siguiente línea un nuevo tipo de columna para tablas: ancho personalizado y contenido centrado
\newcolumntype{P}[1]{>{\centering\arraybackslash}p{#1}}
\DeclareCaptionFormat{upper}{#1#2\uppercase{#3}\par}

% Diseño de tabla para ingeniería
\captionsetup[table]{
	format=upper,
	justification=centering,
	labelsep=period,
	width=.75\linewidth,
	labelfont=small,
	font=small,
}

%Diseño de tabla para ciencias sociales y humanidades
%\captionsetup[table]{
%	justification=raggedright,
%	labelsep=period,
%	labelfont=small,
%	singlelinecheck=false,
%	font={small,bf}
%}

\usepackage{tabularx}

% DISEÑO DE FIGURAS. Puedes elegir entre el estilo para ingeniería o para ciencias sociales y humanidades. Por defecto, está activado el estilo de ingeniería. Si deseas utilizar el otro, comenta las líneas del diseño de ingeniería y descomenta las del diseño de ciencias sociales y humanidades
\usepackage{graphicx}
\graphicspath{{imagenes/}} %ruta a la carpeta de imágenes

% Diagramas conceptuales (TikZ) — pipeline, arquitectura de agentes
\usepackage{tikz}
\usetikzlibrary{positioning, arrows.meta, shapes.geometric, fit, backgrounds, calc}

% Diseño de figuras para ingeniería
\captionsetup[figure]{
	format=hang,
	name=Fig.,
	singlelinecheck=off,
	labelsep=period,
	labelfont=small,
	font=small		
}

% Diseño de figuras para ciencias sociales y humanidades
%\captionsetup[figure]{
%	format=hang,
%	name=Figure,
%	singlelinecheck=off,
%	labelsep=period,
%	labelfont=small,
%	font=small		
%}


% NOTAS A PIE DE PÁGINA
\usepackage{chngcntr} %para numeración contínua de las notas al pie
\counterwithout{footnote}{chapter}

% LISTADOS DE CÓDIGO
% soporte y estilo para listados de código. Más información en https://es.wikibooks.org/wiki/Manual_de_LaTeX/Listados_de_código/Listados_con_listings
\usepackage{listings}

% definimos un estilo de listings
\lstdefinestyle{estilo}{ frame=Ltb,
	framerule=0pt,
	aboveskip=0.5cm,
	framextopmargin=3pt,
	framexbottommargin=3pt,
	framexleftmargin=0.4cm,
	framesep=0pt,
	rulesep=.4pt,
	backgroundcolor=\color{gray97},
	rulesepcolor=\color{black},
	%
	basicstyle=\ttfamily\footnotesize,
	keywordstyle=\bfseries,
	stringstyle=\ttfamily,
	showstringspaces = false,
	commentstyle=\color{gray45},     
	%
	numbers=left,
	numbersep=15pt,
	numberstyle=\tiny,
	numberfirstline = false,
	breaklines=true,
	xleftmargin=\parindent
}

\captionsetup[lstlisting]{font=small, labelsep=period}
% fijamos el estilo a utilizar 
\lstset{style=estilo}
\renewcommand{\lstlistingname}{\uppercase{Código}}


%BIBLIOGRAFÍA - PUEDES ELEGIR ENTRE ESTILO IEEE O APA. POR DEFECTO ESTÁ CONFIGURADO IEEE. SI DESEAS USAR APA, COMENTA LAS LÍNEA DE IEEE Y DESCOMENTA LAS DE APA. Si haces cambios en la configuración de la bibliografía y no obtienes los resultados esperados, es recomendable limpiar los archivos auxiliares y volver a compilar en este orden: COMPILAR-BIBLIOGRAFIA-COMPILAR
% Tienes más información sobre cómo generar bibliografía en http://tex.stackexchange.com/questions/154751/biblatex-with-biber-configuring-my-editor-to-avoid-undefined-citations , https://es.sharelatex.com/learn/Bibliography_management_in_LaTeX y en http://www.ctan.org/tex-archive/macros/latex/exptl/biblatex-contrib
% También te recomendamos consultar la guía temática de la Biblioteca sobre citas bibliográficas: http://uc3m.libguides.com/guias_tematicas/citas_bibliograficas/inicio

% CONFIGURACIÓN PARA LA BIBLIOGRAFÍA IEEE
\usepackage[backend=biber, style=ieee, isbn=false,sortcites, maxbibnames=5, minbibnames=1]{biblatex} % Configuración para el estilo de citas de IEEE, recomendado para el área de ingeniería. "maxbibnames" indica que a partir de 5 autores trunque la lista el primero (minbibnames) y añada "et al." tal y como se utiliza en el estilo IEEE.

%CONFIGURACIÓN PARA LA BIBLIOGRAFÍA APA
%\usepackage[style=apa, backend=biber, natbib=true, hyperref=true, uniquelist=false, sortcites]{biblatex}
%\DeclareLanguageMapping{spanish}{spanish-apa}

\addbibresource{bibliografia/bibliografia.bib} % llama al archivo bibliografia.bib que utilizamos de ejemplo

\usepackage{pdflscape}



%-------------
%	DOCUMENTO
%-------------

\begin{document}
\pagenumbering{roman} % Se utilizan cifras romanas en la numeración de las páginas previas al cuerpo del trabajo
	
%----------
%	PORTADA
%----------	
\begin{titlepage}
	\begin{sffamily}
	\color{azulUC3M}
	\begin{center}
		% Logo Universidad - Tamaño ajustado
        
\includegraphics[width=6cm]{uc3m-logo-171x60.png} 
        
        \vspace{0.8cm} % Espacio entre logos
        
        % Logo RTVE Impulsa - Mucho más pequeño para que no domine la página
        
\includegraphics[width=5.5cm]{impulsa_rtve.jpeg}
        
        \vspace{1.5cm}

        {\Huge ``Postedición Asistida por IA: \\Detección de Sesgos de Género en la Representación Mediática de la Mujer''}\\
		 	\vspace*{0.5cm}
	 		\rule{10.5cm}{0.1mm}\\
			\vspace*{0.9cm}
			{\LARGE Jorge Garcelán Gómez}\\ 
			\vspace*{1cm}
            
		\begin{Large}
			Máster en Inteligencia Artificial Aplicada\\			
			2025-2026 \\
			\vspace{2cm}		
			\textsl{Trabajo Final de Máster}
			\bigskip
			
		\end{Large}
		 	
		\begin{Large}
			Carmen Peláez Moreno\\
			Teresa Onorati\\
			Madrid, Julio 2026\\
		\end{Large}
	\end{center}
	\vfill
	\color{black}
	% si nuestro trabajo se va a publicar con una licencia Creative Commons, incluir estas líneas. Es la opción recomendada.
	
\includegraphics[width=4.2cm]{imagenes/creativecommons.png}\\  %incluimos el logotipo de creativecommons
	This work is licensed under Creative Commons \textbf{Attribution – Non Commercial – Non Derivatives}
	\end{sffamily}
\end{titlepage}

\newpage %página en blanco o de cortesía
\thispagestyle{empty}
\mbox{}

%----------
%	RESUMEN Y PALABRAS CLAVE
%----------	
\renewcommand\abstractname{\large\uppercase{Resumen}}
\begin{abstract}
\thispagestyle{plain}
\setcounter{page}{3}
	
	% ESCRIBIR EL RESUMEN AQUÍ
    \color{red}
    
    La transformación digital del periodismo y los medios de comunicación ha impulsado la integración de herramientas de Inteligencia Artificial (IA) en las redacciones, una oportunidad que, no obstante, conlleva riesgos significativos si no se gestiona bajo una perspectiva ética y socialmente responsable. En particular, la representación de la mujer en las noticias sigue padeciendo de desequilibrios estructurales, tanto cuantitativos, con una presencia minoritaria de mujeres como fuentes expertas y protagonistas \cite{unwomen2025gmmp}, como cualitativos, a través de discursos que perpetúan estereotipos de género (Kotek, 2023). Esta problemática se puede agravar con el uso de Grandes Modelos de Lenguaje, que al haber sido entrenados con inmensos volúmenes de datos, pueden tender a alucinar y amplificar sesgos preexistentes (Navigli 2023). Ante este escenario, las soluciones actuales de detección de sesgos resultan insuficientes (Lai et al. 2021): los métodos manuales carecen de escalabilidad para el ritmo de una redacción moderna, mientras que los sistemas totalmente automatizados a menudo fallan en captar los matices sutiles del lenguaje periodístico.
    
    En respuesta a este desafío, el presente proyecto propone una solución tecnológica híbrida bajo el paradigma human-in-the-loop (“persona en el bucle”). El objetivo principal es investigar y desarrollar una herramienta de apoyo a la postedición que permita a los periodistas identificar y corregir sesgos de género antes de la publicación, sin sustituir su criterio editorial. Para lograrlo, se plantea una arquitectura de software innovadora basada en sistemas multi-agente (Li et al., 2023) y en el uso de técnicas avanzadas de prompt engineering (“ingeniería de instrucciones”), específicamente la metodología Chain-of-Thought (“Cadena de Pensamiento”) (Wei et al., 2022), que obliga a los modelos a razonar y explicitar los motivos por los que consideran un fragmento sesgado, aumentando así la transparencia y la confianza del usuario final. Además, para mitigar las alucinaciones propias de los LLM, se diseñará un sistema de Generación Aumentada por Recuperación (RAG). Esto permitirá al sistema contrastar la información generada con una base de conocimiento validada y experta, asegurando que las sugerencias de corrección no solo sean gramaticalmente correctas, sino periodísticamente rigurosas.

    Finalmente, la robustez de la propuesta reside en su metodología de validación. Para garantizar la eficacia del sistema, se utilizará como ground truth (verdad absoluta) el corpus del proyecto IRIS\_IAMEDIA, un conjunto de datos previamente analizado y etiquetado por expertas en comunicación y género. Este enfoque permitirá entrenar y evaluar los modelos comparando sus diagnósticos directamente con el criterio humano experto, asegurando así que la herramienta se alinee con los estándares éticos requeridos para su transferencia real a los medios de comunicación.

    \color{black}
	
	\textbf{Palabras Clave:}
	% Escribir las palabras clave aquí
     \textit{Sesgos de género, Inteligencia Artificial, Periodismo Computacional, Agentes Autónomos, Grandes Modelos de Lenguaje, Persona en el bucle.}
	
	\vfill
\end{abstract}
	\newpage % página en blanco o de cortesía
	\thispagestyle{empty}
	\mbox{}


%----------
%	DEDICATORIA
%----------	
\chapter*{Dedicatoria}

\setcounter{page}{5}
	
	% ESCRIBIR LA DEDICATORIA AQUÍ	
Llegar hasta aquí ha sido un reto enorme, y ver este TFM terminado me llena de orgullo. Este trabajo está dedicado, antes que a nadie, a mi familia: a mis padres, al yayo y a la yaya, y a mis familiares más cercanos. Gracias por ser ese empujón constante y por ser mi lugar seguro siempre.

A mis amigos, por ser mi vía de escape favorita con nuestras charlas y risas, y por llenar esta etapa de anécdotas y momentos que valen oro.

A nivel profesional, me llevo grandes aprendizajes de esta etapa. Por un lado, quiero dar las gracias a mis compañeros y compañeras de AWS por acompañarme durante mi etapa allí y hacer de esos meses una experiencia fantástica.

\color{red}
Quiero agradecer también a RTVE...
\color{black}

Por otro lado, y para terminar, gracias de corazón a Carmen y Teresa. Haber podido compaginar este máster con mi contrato en el proyecto IRIS ha marcado un antes y un después. Trabajar codo con codo en un entorno de investigación real me ha dado una perspectiva del mundo académico increíble. Gracias por aterrizar mis ideas y por remar conmigo para que este proyecto saliera adelante; trabajar a vuestro lado ha sido un lujo.

		
	\vfill
	
	\newpage % página en blanco o de cortesía
	\thispagestyle{empty}
	\mbox{}
	

%----------
%	ÍNDICES
%----------	

%--
%Índice general
%-
\tableofcontents
\thispagestyle{fancy}

\newpage % página en blanco o de cortesía
\thispagestyle{empty}
\mbox{}

%--
% Índice de figuras. Si no se incluyen, comenta las líneas siguientes
%-
\listoffigures
\thispagestyle{fancy}

\newpage % página en blanco o de cortesía
\thispagestyle{empty}
\mbox{}

%--
% Índice de tablas. Si no se incluyen, comenta las líneas siguientes
%-
\listoftables
\thispagestyle{fancy}

\newpage % página en blanco o de cortesía
\thispagestyle{empty}
\mbox{}


%----------
%	TRABAJO
%----------	
\clearpage
\pagenumbering{arabic} % numeración con múmeros arábigos para el resto de la publicación	
\renewcommand{\footnotesize}{\tiny}  % Changes footnote size to \scriptsize

    %----------
    %	INTRODUCTION
    %----------	

\chapter{Introducción}
\label{introduction}
%----------
%	INTRODUCTION
%----------	

La aplicación de la Inteligencia Artificial para abordar problemáticas sociales críticas encuentra un claro exponente en el proyecto Bindi, impulsado por el grupo multidisciplinar UC3M4Safety de la Universidad Carlos III de Madrid (UC3M). Este equipo ha desarrollado una tecnología pionera basada en aprendizaje automático capaz de detectar indicios de violencia de género a partir de biomarcadores paralingüísticos de la voz. Sin embargo, a pesar de que este avance científico fue liderado y materializado por investigadoras, destacando la labor de Carmen Peláez Moreno, Celia López Ongil, Clara Sáinz de Baranda Andújar, Emma Reyner Fuentes y Esther Rituerto González, una parte notable de la cobertura mediática difundió el hallazgo atribuyendo el mérito a un equipo de «investigadores», como en la Imagen \ref{investigadores_bindi}. El uso automático de este masculino genérico invisibilizó la autoría y el liderazgo femenino en una tecnología diseñada específicamente para la protección de las mujeres, evidenciando un sesgo de representación estructural que sigue profundamente arraigado en las rutinas de redacción periodística actual.


\begin{figure}[H]
	\ffigbox[\FBwidth]
	{\caption{Ejemplo de noticia publicada haciendo mención a «investigadores». }\label{investigadores_bindi}} 
 % [Aquí ponemos como queremos que aparezca en el índice, sin la referencia]{Aquí como queremos que aparezca en el documento, con la referencia}
	{
\includegraphics[scale=0.3]{Plantilla_TFG_ingles_2019/imagenes/bindi_investigadores.png}}
\end{figure}


En este contexto, la transformación digital de los medios de comunicación ha acelerado la integración de herramientas de Inteligencia Artificial en las redacciones. Esta oportunidad, no obstante, conlleva riesgos significativos si no se gestiona bajo una perspectiva ética. Como ilustra el caso de Bindi, la representación de la mujer en las noticias padece desequilibrios estructurales, tanto cuantitativos, con una presencia minoritaria como fuentes expertas \cite{unwomen2025gmmp}, como cualitativos, a través de discursos que perpetúan estereotipos \cite{kotek2023}. Esta problemática amenaza con agravarse con la adopción de Grandes Modelos de Lenguaje (LLM), los cuales, al haber sido entrenados con inmensos volúmenes de textos históricos, tienden a heredar, alucinar y amplificar estos sesgos preexistentes \cite{navigli2023}. Ante este escenario, las soluciones convencionales resultan insuficientes \cite{lai2021}: los métodos de revisión manual carecen de escalabilidad para el ritmo de una redacción moderna, mientras que los sistemas de detección puramente automatizados fallan al captar los matices sutiles del lenguaje periodístico.

En respuesta a este desafío, el presente Trabajo de Fin de Máster propone una solución tecnológica híbrida fundamentada en el paradigma \textit{human-in-the-loop} («persona en el bucle»). El objetivo principal es investigar y desarrollar una herramienta de apoyo a la postedición que permita a los y las profesionales de la información identificar y corregir sesgos de género antes de la publicación, sin sustituir su criterio editorial. Para lograrlo, se plantea una arquitectura basada en sistemas multi-agente \cite{li2023} y en el uso de técnicas avanzadas de \textit{prompt engineering}. Específicamente, se emplea la metodología \textit{Chain-of-Thought} \cite{wei2023}, que obliga a los modelos a razonar y explicitar los motivos por los que consideran un fragmento sesgado, aumentando la transparencia y la confianza del resultado final. Además, para mitigar las alucinaciones propias de los LLM, se integra un sistema de Generación Aumentada por Recuperación (RAG). Esto permite contrastar la información generada con una base de conocimiento validada y experta, asegurando que las sugerencias de corrección sean periodísticamente rigurosas.

Finalmente, la robustez de la propuesta reside en su metodología de validación. Para garantizar la eficacia y fiabilidad del sistema, se utiliza como \textit{ground truth} el corpus del proyecto IRIS\_IAMEDIA, un conjunto de datos previamente analizado y etiquetado por expertas en comunicación y perspectiva de género. Este enfoque permite evaluar los modelos comparando sus diagnósticos directamente con el criterio humano experto, asegurando así que la herramienta cumpla con los estándares éticos requeridos para su transferencia real a la industria de los medios de comunicación.








\section{Objetivos y Motivación}
\label{sec:objectives_motivation}
%----------
%	OBJECTIVES AND MOTIVATION
%----------	

Esta sección está dedicada a la presentación de los objetivos generales y específicos que se pretenden lograr con el desarrollo de este Trabajo Final de Máster. Podemos identificar los siguientes objetivos:

\textbf{OG1.} Investigar las arquitecturas y métodos de diseño, entrenamiento y refinado de los modelos de IA más avanzados, incluídos los LLM y sistemas basados en agentes, para analizar su viabilidad y prestaciones para detectar y localizar variables cuantitativas y cualitativas relacionadas con la representación mediática de las mujeres y susceptibles de ser integradas posteriormente en una herramienta que apoye la postedición de noticias.

\textbf{OE1.} Diseñar la arquitectura conceptual del sistema, definiendo los roles funcionales de los agentes, su relación con los flujos de trabajo periodísticos y los mecanismos de interacción entre los distintos módulos (análisis, verificación y recomendación).

\textbf{OE2.} Desarrollar y evaluar estrategias de prompting y herramientas asociadas (criterios, plantillas, contextos expertos) para la identificación coherente de variables discursivas, estilísticas y de sesgo.

\textbf{OE3.} Implementar técnicas de representación y recuperación de información, incluyendo RAG, búsqueda semántica y medidas de similitud, para mejorar el razonamiento contextual de los modelos y enriquecer sus recomendaciones.

\textbf{OE4.} Comparar distintos modelos de lenguaje y pipelines de procesamiento del lenguaje natural a fin de determinar su capacidad para detectar sesgos explícitos e implícitos y su idoneidad para su integración en un sistema de postedición en el futuro.

\textbf{OE5.} Proponer directrices para la futura integración del sistema en entornos editoriales reales, incluyendo escenarios de uso, requisitos éticos, criterios de explicabilidad y mecanismos de supervisión humana.


\newpage

\section{Paquetes de Trabajo (PT) y Planificación}
\label{sec:workmethodology_planning}
%----------
%	WORK METHODOLOGY AND PLANNING
%----------

Para estructurar el desarrollo de la investigación y garantizar el cumplimiento de los objetivos, el proyecto se ha dividido en seis Paquetes de Trabajo (PT). A continuación, se detallan las tareas principales y los entregables asociados a cada uno de ellos:


\begin{itemize}
    \item \textbf{PT0. Coordinación, comunicación y diseminación}
    \begin{itemize}
        \item Reuniones periódicas con el equipo de IRIS\_IAMEDIA.
        \item Supervisión académica y de RTVE.
        \item Documentación y difusión de avances.
        \item Preparación de publicaciones para congresos, conferencias y revistas científicas de ámbito nacional e internacional.
        \item Participación en eventos de divulgación y jornadas de comunicación.
        \item \textbf{Entregable:} Entrega final y defensa del TFM en RTVE y en UC3M.
    \end{itemize}
    
    \item \textbf{PT1. Diseño de la arquitectura agéntica}
    \begin{itemize}
        \item Definición de las agentes especializadas y sus roles.
        \item Diseño del bucle agéntico de razonamiento y acción, con verificación de evidencias.
        \item Especificación del formato de salida estructurado.
        \item Esquemas de interacción entre agentes, herramientas y niveles de contexto.
    \end{itemize}
    
    \item \textbf{PT2. Diseño y evaluación de Agent Skills y prompts}
    \begin{itemize}
        \item Formalización de la metodología experta como \textit{skills} cargables.
        \item Creación de plantillas y estrategias de \textit{prompting}.
        \item Experimentos y pruebas de consistencia, con ajustes iterativos.
    \end{itemize}
    
    \item \textbf{PT3. Representaciones semánticas y retrieval}
    \begin{itemize}
        \item Implementación de RAG.
        \item Índices de similitud.
        \item Evaluación de precisión contextual.
    \end{itemize}
    
    \item \textbf{PT4. Evaluación comparada de modelos}
    \begin{itemize}
        \item Pruebas con modelos comerciales por API y modelos locales de pesos abiertos.
        \item Métricas de rendimiento y de acuerdo.
        \item Identificación de fortalezas y limitaciones.
    \end{itemize}

    \newpage 
    
    \item \textbf{PT5. Prueba de concepto, conclusiones y transferencia}
    \begin{itemize}
        \item Integración con el proyecto IRIS\_IAMEDIA.
        \item Directrices de integración editorial.
        \item Conclusiones y recomendaciones finales.
    \end{itemize}
\end{itemize}



\begin{landscape}
\begin{figure}[H]
	\ffigbox[\FBwidth]
	{\caption[Cronograma de trabajo (diagrama de Gantt)]
 {Cronograma de trabajo (diagrama de Gantt)}}
 % [Aquí ponemos como queremos que aparezca en el índice, sin la referencia]{Aquí como queremos que aparezca en el documento, con la referencia}
	{\includegraphics[scale=0.6]{Plantilla_TFG_ingles_2019/imagenes/cronograma_iris_blanco.pdf}}
\end{figure}
\end{landscape}

\newpage







\section{Estructura}
\label{sec:thesis_structure}
%----------
%	THESIS STRUCTURE
%----------	

La investigación se estructura de la siguiente manera en capítulos y secciones a lo largo de todo el documento:

\begin{table}[H]
    \ttabbox[\FBwidth]{
        \caption{Estructura de la Investigación}
        \label{table:thesis_structure} % Unique label for the table
    }{
        \begin{tabular}{p{3cm} p{11cm}}
            \hline
            \rowcolor{gray!30} % Adding a gray color to the header row
            \textbf{Capítulo / Anexo} & \textbf{Descripción} \\ \hline
            Capítulo 2: Estado del Arte & Revisión de la literatura sobre la representación mediática de la mujer, los sesgos de género y los métodos de detección basados en IA y procesamiento de lenguaje natural, así como los modelos de lenguaje, las agentes y las habilidades agénticas sobre los que se construye la solución, y la justificación de la propuesta. \\ \hline
            
            Capítulo 3: Metodología & Detalla el marco metodológico: el análisis experto con perspectiva de género, el corpus del proyecto IRIS\_IAMEDIA, la selección de las cinco variables de lenguaje y los modelos de lenguaje empleados. \\ \hline

            Capítulo 4: Implementación & Describe el desarrollo técnico de la solución: la arquitectura de agentes especializadas, el empaquetado de la metodología como Agent Skills, la recuperación de información (RAG) y el protocolo de evaluación, en línea con los objetivos de la Sección~\ref{sec:objectives_motivation}. \\ \hline

            Capítulo 5: Resultados y Discusión & Presenta los resultados de la evaluación variable a variable, la discusión de sus causas, repartidas entre el límite de la tarea y el de los modelos, y la aplicabilidad editorial mediante la integración en la herramienta de IRIS. \\ \hline
            
            Capítulo 6: Conclusiones y Líneas Futuras & Sintetiza los hallazgos principales de la investigación, reflexiona sobre el cumplimiento de los objetivos y propone futuras vías de desarrollo y aplicación. \\ \hline
            
            Anexo A: Presupuesto & Desglosa los costes humanos, tecnológicos y de licencias de \textit{software} necesarios para la viabilidad y ejecución de este proyecto. \\ \hline

    
            Anexo B: Código & Proporciona el acceso al repositorio digital que aloja el código fuente y la documentación técnica del sistema desarrollado. \\ \hline


            Anexo C: Cumplimiento del AI Act & Evalúa y justifica la exención del proyecto ante las normativas más restrictivas de la Ley de IA de la Unión Europea, por su carácter de investigación académica. \\ \hline
            Anexo D: Declaración de uso de IAG & Recoge la declaración de uso de Inteligencia Artificial Generativa cumplimentada conforme al modelo de la Universidad, con el uso técnico realizado y una reflexión sobre su utilidad. \\ \hline
        \end{tabular}
    }
\end{table}

\newpage



\section{Marco Regulatorio}
\label{sec:reg_framework}
%----------
%	REGULATORY FRAMEWORK
%----------	

En esta sección se detallan las bases normativas y los principios éticos que enmarcan el desarrollo de la presente investigación.

En primer lugar, dado que el análisis de textos periodísticos conlleva el procesamiento de información que puede identificar a personas físicas (firmas de periodistas, nombres de fuentes expertas o figuras públicas), el proyecto se supedita estrictamente a las directrices del Reglamento General de Protección de Datos (RGPD)\footnote{Reglamento (UE) 2016/679 del Parlamento Europeo y del Consejo. Disponible en: \url{https://eur-lex.europa.eu/legal-content/ES/TXT/?uri=CELEX\%3A32016R0679} [Último acceso: 15/03/2026]}. La metodología garantiza que el tratamiento de estos datos se realiza con un fin analítico centrado en detectar patrones lingüísticos y sesgos de género, eludiendo cualquier tipo de perfilado personal de la información extraída.

En el ámbito tecnológico, el uso de los LLM y de sistemas basados en agentes obliga a contextualizar este TFM dentro del Reglamento de Inteligencia Artificial de la Unión Europea (AI Act)\footnote{Reglamento (UE) 2024/1689 del Parlamento Europeo y del Consejo, de 13 de junio de 2024, por el que se establecen normas armonizadas en materia de inteligencia artificial. Disponible en: \url{https://eur-lex.europa.eu/legal-content/ES/TXT/?uri=CELEX\%3A32024R1689} [Último acceso: 13/09/2026]}. Al tratarse de un proyecto de investigación científica y desarrollo académico, la herramienta queda fuera del ámbito de las obligaciones aplicables a los sistemas de alto riesgo, en virtud de la exclusión que el propio Reglamento prevé para los sistemas desarrollados y puestos en servicio con el único fin de la investigación científica (artículo 2, apartado 6), tal y como se detalla en la evaluación de cumplimiento adjunta en el Anexo~\ref{appendix_ai_act}. No obstante, cabe destacar que el diseño metodológico de este trabajo, basado en el enfoque \textit{human-in-the-loop}, se alinea de forma proactiva con los principios de la normativa europea: exigencia de supervisión humana, transparencia explicativa y mitigación de sesgos discriminatorios.

En lo relativo a la obtención del corpus de estudio, este trabajo se nutre de la recopilación de artículos periodísticos realizada en el marco del proyecto IRIS\_IAMEDIA e InfoIA. Toda la extracción previa de estos datos se llevó a cabo respetando rigurosamente los términos y condiciones de servicio de las cabeceras digitales analizadas. Al tratarse de una investigación académica impulsada por la Universidad Carlos III de Madrid, el procesamiento original de los textos se amparó en las excepciones relativas a la minería de textos y datos (\textit{Text and Data Mining}, TDM por sus siglas en inglés) para fines de investigación científica, heredando este Trabajo de Fin de Máster dichas garantías legales y éticas.

Por último, el desarrollo del código de la investigación se ha realizado en Python con librerías de código abierto. En cuanto a los modelos, se han combinado servicios comerciales accesibles por API con un modelo local de pesos abiertos (\texttt{gemma4:e4b}). La publicación del código y la documentación de la metodología buscan garantizar la transparencia y la reproducibilidad del trabajo.

\newpage

\section{Impacto Socio-Económico}
\label{sec:socioeconomic_impact}
%----------
%	SOCIOECONOMIC IMPACT
%----------	

Esta sección explora cómo el presente TFM contribuye a diversas esferas de la sociedad y la economía, subrayando tanto los beneficios potenciales de su aplicación como los retos inherentes a la automatización.

En el plano social, fomentar narrativas periodísticas libres de sesgos de género contribuye directamente a la construcción de un imaginario colectivo más justo e igualitario. Al facilitar la detección de patrones de desigualdad, esta investigación ayuda a los medios a ofrecer una visión más equilibrada de la realidad. Este cambio positivo no solo empodera a las mujeres ofreciendo referentes diversos, sino que también mejora la calidad democrática de la información, fomentando una audiencia más crítica y mejor informada.

Desde una perspectiva económica e industrial, la herramienta propuesta resulta de gran valor para los medios de comunicación y las corporaciones públicas. Tradicionalmente, las auditorías de calidad y perspectiva de género exigen un ingente esfuerzo manual. Al automatizar este primer filtro mediante agentes especializados y LLM, las redacciones pueden procesar grandes volúmenes de texto de manera eficiente, optimizando los flujos de trabajo editoriales y reduciendo significativamente los costes operativos.

Más allá del entorno periodístico y académico, los resultados de este trabajo pueden ser utilizados por instituciones públicas y observatorios de igualdad para diseñar políticas, normativas e iniciativas orientadas a garantizar la igualdad en los medios de comunicación.

El estudio exhaustivo del presupuesto económico empleado para la viabilidad y desarrollo de este proyecto se encuentra detallado en el Anexo \ref{appendix_budget}.

\newpage

\section{Riesgos y Medidas de Mitigación}
\label{sec:risks}
%----------
%	RIESGOS Y MEDIDAS DE MITIGACIÓN
%----------	
La Tabla~\ref{table:risks} recoge los principales riesgos identificados al planificar el trabajo, asociados a los Paquetes de Trabajo (PT) descritos en la Sección~\ref{sec:workmethodology_planning}, junto con su valoración y las medidas de mitigación previstas entonces. Como esta memoria se redacta una vez concluido el trabajo, al análisis inicial le sigue un balance de lo que finalmente ocurrió con cada riesgo y de en qué medida las respuestas previstas resultaron aplicables.


\begin{table}[H]
    \small % Font size
    \ttabbox[\FBwidth]{
        \caption{Riesgos y medidas de mitigación}
        \label{table:risks} % Etiqueta única para la tabla
    }{
        \begin{tabular}{p{4cm} p{2cm} p{2cm} p{1cm} p{4cm}}
            \hline
            \rowcolor{gray!30} % Adding a gray color to the header row
            \textbf{Riesgo} & \textbf{Verosimilitud} & \textbf{Severidad} & \textbf{PTs} & \textbf{Medidas de Mitigación} \\ \hline
            
            \textbf{R1.} Limitaciones de recursos provenientes de IRIS\_IAMEDIA (anotaciones, validación experta, corpus ampliado) & Baja, Media & Media & Todos & Uso del corpus ya validado de IRIS\_IAMEDIA, anotación complementaria propia, generación de ejemplos sintéticos con LLM, retroalimentación puntual en lugar de continua. \\ \hline

            
            \textbf{R2.} Variabilidad y falta de consistencia en los LLM al detectar sesgos & Media & Alta & PT2, PT3, PT4 & Uso de prompting estructurado, definición de casos de test, comparación entre modelos, agente verificador y validación con corpus experto IRIS\_IAMEDIA. \\ \hline

            
            \textbf{R3.} Dificultad para capturar variables cualitativas complejas (androcentrismo, infantilización, etc.) & Media & Media, Alta & PT2, PT4 & Ajuste iterativo de prompts, RAG para aportar contexto, supervisión \textit{human-in-the-loop}, casos anotados por expertas. \\ \hline

            
            \textbf{R4.} Desalineación entre los diagnósticos del sistema y las prácticas y necesidades reales de las redacciones & Media & Media & PT1, PT5 &  Intensificar las consultas con expertas, incorporación de casos de uso reales, test de usabilidad y ajustes iterativos. \\ \hline

            
            \textbf{R5.} Dependencia de herramientas externas o carga académica que afecte al calendario de ejecución & Media & Media & Todos & Planificación semanal, priorización de entregables y alternativas técnicas si fallan las API o servicios externos. \\ \hline

        \end{tabular}
    }
\end{table}

Concluido el trabajo, el balance es el siguiente. Tres de los cinco riesgos se
materializaron, y en ellos las medidas previstas resultaron aplicables, aunque no
siempre suficientes.

El \textbf{R1} se materializó de forma parcial. El corpus anotado estuvo disponible y
bastó para evaluar el plano cuantitativo del sistema, pero no incluía doble
codificación ni marcado experto sobre el texto. Esa carencia impidió calcular la
fiabilidad intercodificadora y validar el plano cualitativo
(Sección~\ref{subsec:limitacion_fiabilidad}). Las medidas previstas de anotación
complementaria y de generación de ejemplos sintéticos no llegaron a emplearse. En su
lugar, la limitación se convirtió en objeto de análisis, ya que el estudio de
sensibilidad por equipo de anotación permitió acotar de forma indirecta el techo de
acuerdo de la tarea (Sección~\ref{subsec:iris_errores}).

El \textbf{R2} se materializó con claridad y acabó siendo uno de los hallazgos del
trabajo, pues los cuatro modelos evaluados mostraron comportamientos marcadamente
distintos entre sí. Se aplicaron las medidas previstas de salida estructurada,
verificación automática de las evidencias, comparación sistemática entre modelos y
validación contra el corpus experto. A ellas se sumó una medida no contemplada
inicialmente, la desactivación del razonamiento explícito en los modelos que lo
admiten, para que la comparación fuese justa
(Sección~\ref{sec:genai_models_iris}).

El \textbf{R3} se materializó y constituye el resultado central del trabajo. La
recuperación aumentada (RAG), prevista como medida de mitigación, se implementó y se
evaluó de forma aislada, con el resultado de que no compensa en este dominio
(Sección~\ref{subsec:iris_descomposicion}). Poder descartarla con evidencia es en sí
mismo una aportación. La supervisión \textit{human-in-the-loop}, en cambio, sí se
incorporó y quedó materializada en la integración del sistema en la herramienta
del proyecto (Sección~\ref{poc}).

El \textbf{R4} no llegó a evaluarse de forma sistemática. La integración del sistema
en la herramienta web del proyecto (Sección~\ref{poc}) y las directrices de uso
editorial que cierran el trabajo (Sección~\ref{sec:conclusions}) responden a la
medida prevista, pero la validación con redacciones reales queda como línea de
trabajo futura (Sección~\ref{sec:futurework_usecases}).

El \textbf{R5} se materializó en dos frentes concretos, y en ambos funcionó la
contingencia prevista. Por un lado, OpenAI anunció la retirada de modelos de su
interfaz, lo que motivó incorporar un modelo más reciente para no depender de uno
solo (Sección~\ref{sec:genai_models_iris}). Por otro, las ejecuciones sobre la granja
de GPU sufrieron interrupciones, para lo que se implementó una ejecución reanudable y
un mecanismo de corte que evita registrar los fallos de infraestructura como
predicciones (Sección~\ref{eval}).


\newpage

\section{Contexto de la Investigación}
\label{sec:research_context}
%----------
%	RESEARCH CONTEXT
%----------	

Este Trabajo de Fin de Máster se desarrolla bajo un doble contexto académico e institucional. Por un lado, se desarrolla en contacto con el proyecto IRIS\_IAMEDIA~\cite{irisiamedia}, del que toma el corpus y la metodología, sin formar parte de él. Se trata de un proyecto interdisciplinar en el que colaboran expertas en comunicación e ingenieras e ingenieros expertos en el ámbito de la Inteligencia Artificial~\cite{avi2026}. Aunque el TFM se nutre del ecosistema científico y tecnológico generado por IRIS\_IAMEDIA, ninguno de sus resultados forma parte de los objetivos, entregables ni compromisos oficiales del proyecto, cuya agenda de investigación es independiente. No obstante, la participación de las investigadoras principales (IP) y colaboradoras de IRIS\_IAMEDIA ha permitido disponer de un marco metodológico consolidado, así como de recursos, corpus experimentales y líneas de trabajo en el análisis de sesgos de género que han inspirado y orientado esta propuesta.

Por otro lado, esta investigación se desarrolla tras haber obtenido el primer puesto en la resolución de la VIII Convocatoria Impulsa Visión RTVE de ayudas a la investigación para Trabajos Fin de Máster\footnote{Resolución de la VIII Convocatoria Impulsa Visión RTVE Ayudas a la Investigación para TFM. Disponible en: \url{https://www.rtve.es/rtve/20260304/resultados-viii-impulsa-vision-rtve-ayudas-a-investigacion-tfm-inteligencia-artificial-periodismo-convocatoria/16964519.shtml} [Último acceso: 09/03/2026].}, dentro de la línea de \textit{Inteligencia artificial y periodismo}. Este reconocimiento, otorgado durante la fase de evaluación de la idea y propuesta del proyecto, avala el interés y la necesidad del sector periodístico por contar con este tipo de soluciones. La convergencia de este marco institucional con el respaldo del proyecto IRIS\_IAMEDIA proporciona el escenario idóneo para explorar herramientas tecnológicas aplicables a los desafíos de las redacciones actuales.




	%----------
    %	STATE-OF-THE-ART
    %----------	

\chapter{Estado del Arte}
\label{state-of-the-art}
Este capítulo establece el marco teórico y tecnológico sobre el que se fundamenta esta investigación. En primer lugar, se analiza la evolución histórica y actual de la representación de las mujeres en los medios de comunicación, evidenciando la persistencia de sesgos de género. Posteriormente, se realiza una revisión exhaustiva de las metodologías y herramientas desarrolladas para la detección de dichos sesgos, abarcando desde el análisis manual hasta las aproximaciones más recientes basadas en Inteligencia Artificial.

\section{Representación de las Mujeres en los Medios}
\label{sec:rep_women_media}
%----------
%	REPRESENTACIÓN DE LAS MUJERES EN LOS MEDIOS
%----------	

La desigualdad en la representación mediática no es un fenómeno reciente, sino un reflejo histórico de asimetrías sociales profundas que se manifiesta de forma evidente en la cobertura de crisis globales. Un caso paradigmático que ilustra la complejidad de este sesgo editorial es la cobertura internacional del tsunami del Océano Índico en 2004. Tal y como denunció la galardonada periodista de investigación y experta en género Deepa Kandaswamy en Women's eNews (2005) \cite{kandaswamy2005}, la respuesta mediática incurrió en una doble negligencia hacia las mujeres: la hipervisibilización sensacionalista y la posterior invisibilización estructural.

Según documentó Oxfam Internacional \cite{oxfam2005}, en diversas zonas afectadas de Indonesia, India y Sri Lanka fallecieron entre tres y cuatro mujeres por cada hombre. Esta drástica brecha de mortalidad fue el resultado directo de roles sociales fuertemente arraigados: las mujeres fallecieron en mayor proporción por encontrarse en sus hogares cuidando de menores o ancianos, por estar esperando en la costa el regreso de los barcos pesqueros, o por carecer de habilidades de supervivencia reservadas culturalmente a los hombres en esas regiones, como trepar árboles o nadar. Como consecuencia del desastre, las supervivientes se enfrentaron a vulnerabilidades extremas derivadas de esta alteración demográfica, incluyendo agresiones sexuales en campos de desplazados, matrimonios forzados de menores y una falta de atención institucional a la pérdida de sus medios de vida.

A pesar de la gravedad de los datos, la prensa \textit{mainstream} ignoró sistemáticamente el informe de Oxfam. Más revelador aún fue el rechazo dentro de la propia industria periodística: cuando se cuestionó esta falta de cobertura en foros internacionales de prensa, las reclamaciones fueron minimizadas por otros profesionales bajo acusaciones de tener una ``mentalidad de gueto femenino'' o un exceso de ``sensibilidad de género'' \cite{kandaswamy2005}. Este episodio histórico demuestra cómo la falta de una mirada crítica en las redacciones no solo genera narrativas incompletas, sino que encubre negligencias estructurales y perpetúa la revictimización.



Lejos de haberse corregido de forma natural, más de dos décadas después aquel desequilibrio estructural apenas ha cambiado. Según los hallazgos del Proyecto de Monitoreo Global de Medios (\textit{Global Media Monitoring Project}, GMMP por sus siglas en inglés) \cite{unwomen2025gmmp} de 2025, el progreso hacia la igualdad de género en las noticias se encuentra estancado, ya que las mujeres solo representan el 26\% de las personas que aparecen, hablan o sobre las que se lee en las noticias impresas y audiovisuales a nivel mundial. Las mujeres continúan siendo una minoría cuando se trata de figurar como protagonistas de la información, constituyendo únicamente el 24\% de los sujetos de las noticias. La brecha se mantiene igual de grave cuando se buscan voces de autoridad, representando apenas el 23\% de las personas consultadas como expertas o portavoces.

Más recientemente, el proyecto de investigación InfoIA publicado en febrero de 2026 por Clara Sáinz de Baranda y Heidy Ríos \cite{sainzdebaranda2026}, que analiza las informaciones sobre IA en los medios digitales españoles, advierte que este ecosistema informativo consolida un ``techo de cristal mediático'' al ser un espacio fuertemente masculinizado. Los datos de este informe revelan que los hombres protagonizan el 41,60\% de las piezas periodísticas frente a un escaso 5,33\% de protagonismo exclusivamente femenino, y acaparan el 61,14\% de las fuentes de autoridad consultadas, relegando a las mujeres al 21,45\%. 

Ahora bien, ambas caras del problema importan. A la infrarrepresentación cuantitativa se suma un sesgo cualitativo que opera en la propia forma del lenguaje, en cómo se nombra y se describe a las mujeres. Se manifiesta, por ejemplo, en el masculino genérico que invisibiliza la autoría femenina (un equipo liderado por investigadoras reducido a ``un equipo de investigadores'', como en el caso de Bindi), en la denominación de las profesionales por sus relaciones familiares en lugar de por su cargo (``la mujer de'' o ``la hija de'' en vez de su nombre y su función) o en el tratamiento asimétrico que describe a una mujer por su aspecto, su edad o su vestimenta mientras presenta a un hombre por su puesto o sus logros.

Este trabajo aborda esa dimensión lingüística y lo hace en sus dos planos. En el plano cuantitativo, detecta la presencia o ausencia de cada fenómeno en cada pieza, lo que permite medir su prevalencia y validarla frente al criterio experto, por ejemplo en qué proporción de las noticias aparece el masculino genérico. En el plano cualitativo, no se queda en la etiqueta: localiza en el texto el fragmento concreto que motiva el veredicto y aporta una explicación, de modo que la analista pueda auditar cada decisión en lugar de aceptarla a ciegas. Ambos planos se operacionalizan en el Capítulo~\ref{sec:methodology}. Detectar estos fenómenos de forma sistemática ha motivado toda una línea de métodos, que se revisa a continuación.






\section{Métodos para la Detección de Sesgos de Género}
\label{sec:det_sesgos}
%----------
%	MÉTODOS DE DETECCIÓN DE SESGOS DE GÉNERO
%----------	

La detección y análisis de sesgos de género en el discurso mediático han experimentado una profunda transformación metodológica en la última década, impulsada por la necesidad de procesar volúmenes de información cada vez mayores. Conviven en ella dos grandes enfoques: el tradicional, basado en la codificación humana, y el automatizado, apoyado en técnicas de Inteligencia Artificial.


\subsection{Enfoques Tradicionales y Análisis Manual}

El análisis de la representación mediática y la publicación de los informes que se derivan de él descansan, todavía hoy, en un procedimiento manual. Los equipos de investigación y las expertas en género elaboran en primer lugar una metodología exhaustiva y detallada, materializada a menudo en libros de códigos extensos, como el que describe la tesis doctoral de Clara Sáinz de Baranda \cite{sainzdebaranda2013}. Después llega el trabajo de campo, que exige recorrer periódico a periódico y noticia a noticia, codificando cada pieza para el análisis cuantitativo, con ayuda de diversas herramientas y aplicaciones, y marcando el texto a mano para el cualitativo.


Este enfoque destaca por su alto nivel de precisión humana y comprensión del contexto sociocultural, siendo el estándar de referencia (\textit{golden standard} o \textit{ground truth}, en inglés) en investigaciones académicas. Como ya se ha mencionado, un ejemplo paradigmático de este rigor es la tesis doctoral de Clara Sáinz de Baranda \cite{sainzdebaranda2013}, que requirió el análisis manual de miles de piezas periodísticas para visibilizar la marginación de las mujeres en la prensa desde 1978 hasta 2010. Del mismo modo, proyectos recientes de gran envergadura como el ya mencionado InfoIA 2026 \cite{sainzdebaranda2026} siguen apoyándose en metodologías de revisión experta para garantizar la fiabilidad de los datos. Sin embargo, estas metodologías presentan limitaciones evidentes en términos de escalabilidad: conllevan una inmensa inversión de tiempo, una altísima carga de trabajo manual y la imposibilidad de auditar el ecosistema mediático en tiempo real.


\subsection{Transición hacia la Automatización e Inteligencia Artificial}
Para superar estas barreras de la escalabilidad manual, en los últimos años han surgido diversas herramientas computacionales y plataformas de monitorización que aplican técnicas de Procesamiento de Lenguaje Natural (PLN) e Inteligencia Artificial para la detección de sesgos y discursos discriminatorios. Estas iniciativas se pueden clasificar en tres grandes grupos según su complejidad técnica:

\begin{itemize}
    \item \textbf{Sistemas basados en diccionarios y glosarios:} Un primer acercamiento a la automatización en las redacciones es el uso de sistemas de validación léxica. Destaca el caso del diario argentino La Nación, que desarrolló la herramienta interna \textit{Genie} y su versión en abierto \textit{Cómo lo digo} \cite{lanacion2024}. Esta aplicación funciona cruzando el texto de la periodista con un glosario predefinido que detecta términos inapropiados o que, por su contexto, podrían perpetuar estereotipos. La herramienta marca estas expresiones y sugiere ``Palabras que se recomienda reemplazar'' o ``Palabras que se recomienda revisar'', actuando como un corrector ortotipográfico con perspectiva de género.

    \item \textbf{Monitorización del discurso:} A nivel macroscópico, existen rastreadores automatizados que auditan la producción periodística de forma masiva. El \textit{GenderGap Tracker}, desarrollado por Informed Opinions \cite{gendergap2018}, analiza diariamente los principales medios de comunicación de Canadá para cuantificar la proporción de voces femeninas frente a masculinas en las noticias. En el ámbito hispanohablante destaca \textit{Hatemedia} \cite{hatemedia2023}, un monitor español de odio en la red. Sin embargo, el hito institucional más relevante en este ámbito es la reciente puesta en marcha de \textit{HODIO} (Huella del Odio y la Polarización) \cite{hodio2026}. Impulsada por el Observatorio Español del Racismo y la Xenofobia (OBERAXE) del Ministerio de Inclusión, Seguridad Social y Migraciones, esta plataforma constituye un esfuerzo pionero desde la administración pública. \textit{HODIO} emplea algoritmos de Inteligencia Artificial para auditar masivamente las principales redes sociales (Instagram, TikTok, X, YouTube y Facebook), para clasificar, cuantificar y monitorizar la prevalencia y evolución de discursos discriminatorios y polarizantes.

    \item \textbf{Modelos avanzados de mitigación y postedición interactiva:} En el punto más avanzado de la investigación se encuentran las herramientas diseñadas no solo para detectar, sino para auditar a los propios algoritmos y ofrecer alternativas de forma interactiva. El proyecto \textit{RECAST} \cite{wright2021} es un sistema interactivo que permite visualizar las predicciones de los modelos de detección de toxicidad y proporciona sugerencias alternativas para el lenguaje marcado como tóxico, facilitando la interpretabilidad y la intervención del usuario sobre el texto.
\end{itemize}

Con todo, la mayoría de estas herramientas operan sobre reglas léxicas o clasificadores cerrados, con una capacidad limitada para razonar sobre el contexto o para justificar sus decisiones. La irrupción de los grandes modelos de lenguaje y de las arquitecturas de agentes abre una vía distinta, que se examina a continuación.


\section{Justificación de la Solución}
\label{sec:justification_solution}
%----------
%	JUSTIFICATION OF THE SOLUTION
%----------	
 
A lo largo de este capítulo, hemos explorado el estado del arte sobre la representación de las mujeres en los medios de comunicación y los sesgos de género presentes en la narrativa periodística. Asimismo, hemos analizado las distintas herramientas de Inteligencia Artificial empleadas para la detección de estos sesgos, comprobando que su integración está reconfigurando los procesos productivos de las redacciones. Sin embargo, esta transformación genera oportunidades a la vez que revela el riesgo de perpetuar desigualdades. Diversos estudios muestran que la representación en la cobertura mediática sigue siendo desigual: más del 75 \% de las fuentes expertas citadas son hombres \cite{unwomen2025gmmp}, un sesgo que los LLM tienden a heredar, asociando nombres femeninos con roles vinculados al cuidado entre tres y seis veces más que los masculinos \cite{kotek2023}.

Ante este panorama, y debido al inmenso volumen de contenido que se publica diariamente, resulta clave contar con metodologías capaces de analizar estos sesgos con rigor y de forma escalable. Sin embargo, como se advierte en \cite{lai2021}, ni los enfoques puramente manuales ni los completamente automatizados pueden abordar por sí solos la complejidad del problema: los primeros son costosos y propensos al error humano, mientras que los segundos pueden generar riesgos de explicabilidad, éticos y legales. Esta limitación ha impulsado el desarrollo de propuestas híbridas que combinan el procesamiento del lenguaje natural, los LLM y la supervisión humana (\textit{human-in-the-loop}), uniendo la potencia del aprendizaje automático con la sensibilidad interpretativa de expertas y expertos en estudios de género y comunicación.

Dentro de este ecosistema de soluciones tecnológicas, y como se ha expuesto en la revisión previa, cabe destacar que, si bien plataformas recientes como \textit{HODIO} suponen un avance analítico sin precedentes para diagnosticar el problema de forma externa y \textit{a posteriori}, su enfoque subraya la necesidad complementaria de herramientas de intervención \textit{a priori}, como la propuesta en este Trabajo de Fin de Máster, que actúen directamente en la propia redacción antes de que el sesgo sea publicado.

Precisamente para dar respuesta a esta necesidad, y en un contexto donde el sector periodístico avanza hacia una integración progresiva de herramientas de IA en los flujos de trabajo editoriales, surge este proyecto, que propone el diseño de un sistema basado en grandes modelos de lenguaje y agentes especializadas que asisten a la analista sin sustituir su criterio. El sistema integra la metodología y el conocimiento experto para identificar y localizar los patrones de sesgo en el texto, y aporta para cada uno las evidencias y una explicación que la profesional revisa y, si procede, corrige, contribuyendo a una comunicación más justa y libre de sesgos.

    %----------
    %	METHODOLOGY
    %----------	

\chapter{Metodología}
\label{sec:methodology}
This Chapter provides a detailed presentation of the different methods, models and procedures that have been used for the elaboration of this thesis.


\section{Análisis experto y perspectiva de género}
\label{sec:experts}
%----------
%	ANÁLISIS EXPERTO Y PERSPECTIVA DE GÉNERO — proyecto IRIS
%----------

El presente trabajo se enmarca en el proyecto IRIS del Instituto Universitario
de Estudios de Género de la Universidad Carlos~III de Madrid, cuyo objetivo es
analizar de forma sistemática la representación de las mujeres en la prensa
española. A diferencia de aproximaciones que abordan la detección de discurso
de odio sobre mensajes breves en redes sociales, el objeto de estudio aquí son
\textbf{piezas periodísticas completas} y el fenómeno analizado no es la
agresión explícita, sino el \textbf{sexismo lingüístico}: los usos del lenguaje
que reproducen desigualdad de forma frecuentemente inadvertida por quien los
emplea.

\subsection{Fundamentación metodológica}

El libro de códigos empleado se apoya en la tradición del análisis de contenido
con perspectiva de género aplicada a medios de comunicación
\cite{sainzdebaranda2013, sainzdebaranda2026}, y operacionaliza criterios
procedentes de guías institucionales de lenguaje no sexista elaboradas por
organismos públicos españoles. Estas guías proporcionan tanto los principios
generales como la casuística concreta a partir de la cual se han derivado las
reglas de decisión de cada variable.

Un concepto resulta central en esta fundamentación: la \textbf{regla de
inversión}. Ante un pasaje candidato, el procedimiento consiste en reescribirlo
intercambiando el género de todas las referencias personales y evaluar si el
resultado se publicaría con naturalidad. Si la versión invertida resulta
extraña, ridícula o irrespetuosa, cabe concluir que la formulación original
incorpora sexismo lingüístico. Este criterio permite distinguir con precisión
entre textos que \emph{hablan sobre} desigualdad —informativos y no
necesariamente sexistas— y textos que \emph{reproducen} la desigualdad en su
forma.

Conviene subrayar que el sexismo lingüístico opera sobre la forma del mensaje y
no sobre su fondo. Una pieza puede presentar un contenido no sexista expresado
mediante formas sexistas, o a la inversa. Esta distinción justifica que el libro
de códigos separe explícitamente las variables de forma de aquellas de
contenido, tal como se detalla en la Sección~\ref{sec:iris_variables}.

\subsection{Procedimiento de anotación}

La anotación del corpus fue realizada por especialistas en comunicación y
estudios de género vinculadas al proyecto, aplicando el libro de códigos sobre
la pieza completa. La unidad de análisis es, por tanto, \textbf{el artículo}: cada
variable recibe un único veredicto por pieza, con independencia del número de
ocurrencias del fenómeno en el texto. Esta decisión, coherente con la práctica
habitual en análisis de contenido, tiene implicaciones relevantes para la
interpretación de los resultados, pues basta una sola ocurrencia para que una
pieza extensa quede codificada como positiva.

\subsection{Limitación: ausencia de doble codificación}
\label{subsec:limitacion_fiabilidad}

Cada noticia del corpus fue codificada por una única anotadora. En consecuencia,
\textbf{no es posible calcular la fiabilidad intercodificadora} del libro de
códigos sobre este material, ni estimar por tanto el techo de acuerdo alcanzable
por un sistema automático.

Esta limitación debe tenerse presente en la interpretación de los resultados del
Capítulo~\ref{results_and_discussion}. Las variables consideradas implican
juicios interpretativos sobre los que dos personas expertas podrían discrepar
legítimamente; en tareas de análisis de contenido con este grado de
subjetividad, los coeficientes de acuerdo entre codificadoras humanas raramente
alcanzan valores próximos a la unidad. Al carecer de esa referencia, cualquier
discrepancia observada entre el sistema y la anotación experta \emph{no puede
atribuirse inequívocamente al error del modelo}: parte de ella podría
corresponder a la variabilidad inherente a la propia tarea de codificación.

Se recomienda, como línea de trabajo futura de bajo coste y alto valor, la doble
codificación de una submuestra del corpus (habitualmente entre el 10\,\% y el
20\,\% en la literatura de análisis de contenido), suficiente para obtener una
estimación defendible de dicha fiabilidad.


\subsection{Fuente de datos}
\label{data_source}
%----------
%	FUENTE DE DATOS — corpus IMIO / IRIS
%----------

El corpus empleado en este trabajo ha sido construido en el marco del proyecto
IRIS y está formado por \textbf{7.115 piezas periodísticas} publicadas en prensa
generalista española. A diferencia de los corpus de redes sociales, habituales
en la detección automática de discurso de odio, se trata de textos extensos
—con una longitud mediana de \textbf{5.881 caracteres}— redactados por
profesionales y sometidos a procesos editoriales, lo que sitúa el fenómeno
analizado en el terreno del sesgo sutil antes que en el de la agresión
explícita.

\subsubsection{Composición temporal y por cabecera}

El corpus abarca dos periodos diferenciados, según se recoge en la
Tabla~\ref{tab:iris-corpus-anios}: una serie histórica que cubre de 2017 a 2021
y una campaña de recogida correspondiente a 2024.

\begin{table}[H]
    \centering
    \begingroup
    \small
    \ttabbox[\FBwidth]{
        \caption{Distribución temporal del corpus IRIS}
        \label{tab:iris-corpus-anios}
    }{
        \begin{tabularx}{\textwidth}{
            >{\centering\arraybackslash}X
            >{\centering\arraybackslash}X
            >{\centering\arraybackslash}X
            >{\centering\arraybackslash}X
            >{\centering\arraybackslash}X
            >{\centering\arraybackslash}X
            >{\columncolor{gray!15}\centering\arraybackslash}X
        }
            \hline
            \textbf{2017} & \textbf{2018} & \textbf{2019} & \textbf{2020} & \textbf{2021} & \textbf{2024} & \textbf{Total} \\
            \hline
            976 & 1.199 & 866 & 527 & 565 & 2.982 & 7.115 \\
            \hline
        \end{tabularx}
    }
    \endgroup
\end{table}

Las piezas proceden de siete cabeceras, con una distribución desigual: la
cabecera más representada aporta 3.117 piezas (43,8\,\% del total) y la menos
representada 63 (0,9\,\%). Esta asimetría, propia de un corpus construido con
criterios de relevancia editorial antes que de balanceo estadístico, no afecta
al diseño experimental de este trabajo —que evalúa el acuerdo pieza a pieza—
pero debe considerarse si se pretendiese extraer conclusiones agregadas sobre
el comportamiento comparado de los medios.

\subsubsection{Obtención del texto completo}

La base de datos original recoge los metadatos de cada pieza (fecha, cabecera,
titular, autoría, sección y sus correspondientes codificaciones), pero no su
cuerpo. Dado que el análisis del sexismo lingüístico requiere necesariamente el
texto íntegro, se desarrolló un proceso de extracción automática que recupera el
contenido a partir de la URL de cada pieza. El resultado se incorpora al corpus
en el campo \texttt{contenido\_articulo}, que constituye la entrada de todos los
experimentos descritos en este trabajo.

\subsubsection{Subconjunto de evaluación}
\label{subsec:iris_subconjunto}

De las 7.115 piezas del corpus, \textbf{1.315 cuentan con anotación experta
simultánea en las cinco variables} objeto de estudio, y son las que constituyen
el conjunto de evaluación. De ellas, 1.313 disponen de texto recuperado con
extensión suficiente para el análisis, cifra sobre la que se calculan todas las
métricas del Capítulo~\ref{results_and_discussion}.

Este subconjunto presenta dos particularidades que condicionan el alcance de las
conclusiones:

\begin{itemize}
    \item \textbf{Restricción temporal.} La totalidad de las piezas anotadas
    corresponde a \textbf{2024}. Los resultados obtenidos no pueden extrapolarse
    sin cautela al periodo 2017--2021, cuyas piezas permanecen sin anotar.

    \item \textbf{Cobertura parcial.} El subconjunto anotado representa el
    18,5\,\% del corpus y procede de cinco de las siete cabeceras, con
    predominio de una de ellas (805 piezas, 61,2\,\% del subconjunto).
\end{itemize}

Las 5.800 piezas restantes carecen de anotación de referencia y, por tanto, no
intervienen en el cálculo de las métricas de acuerdo. Su procesamiento resulta
de interés únicamente para el análisis descriptivo y para la prueba de concepto
editorial, escenarios en los que no se dispone de verdad de referencia contra la
que contrastar.


\subsection{Variables del análisis}
\label{target_variables_defs}
%----------
%	VARIABLES DEL ANÁLISIS — sexismo lingüístico
%----------

El libro de códigos de IRIS define quince variables de sexismo lingüístico,
codificadas como V25 a V39, que se recogen en la
Tabla~\ref{tab:iris-inventario-variables}. Cada una captura un mecanismo
distinto mediante el cual el lenguaje periodístico puede reproducir
desigualdad.

\begin{table}[H]
    \centering
    \begingroup
    \small
    \ttabbox[\FBwidth]{
        \caption{Inventario de variables de sexismo lingüístico del libro de códigos}
        \label{tab:iris-inventario-variables}
    }{
        \begin{tabularx}{\textwidth}{
            >{\centering\arraybackslash}p{1.2cm}
            X
            >{\columncolor{gray!15}\centering\arraybackslash}p{2.2cm}
        }
            \hline
            \textbf{Cód.} & \textbf{Variable} & \textbf{En este trabajo} \\
            \hline
            V25 & \texttt{lenguaje\_sexista} & \textbf{Sí} \\
            V26 & \texttt{masc\_generico} & \textbf{Sí} \\
            V27 & \texttt{hombre\_denominar\_humanidad} & No \\
            V28 & \texttt{uso\_dual\_zorr} & No \\
            V29 & \texttt{uso\_cargo\_mujer} & No \\
            V30 & \texttt{sexismo\_discurso} & \textbf{Sí} \\
            V31 & \texttt{androcentrismo} & No \\
            V32 & \texttt{mencion\_nombre\_investigadora} & No \\
            V33 & \texttt{asimetria\_mujer\_hombre} & \textbf{Sí} \\
            V34 & \texttt{disminutivos\_infantilizacion} & No \\
            V35 & \texttt{denominacion\_sexualizada} & \textbf{Sí} \\
            V36 & \texttt{denominacion\_redundante} & No \\
            V37 & \texttt{denominacion\_dependiente} & No \\
            V38 & \texttt{criterios\_excepcion} & No \\
            V39 & \texttt{comparacion\_mujer\_hombre} & No \\
            \hline
        \end{tabularx}
    }
    \endgroup
\end{table}

\subsubsection{Selección de las cinco variables}

El alcance de este trabajo se restringe a cinco del bloque de Lenguaje (V25, V26, V30, V33 y
V35), seleccionadas por el equipo experto del proyecto atendiendo a su relevancia teórica dentro del marco de análisis. Las diez variables restantes permanecen como
marco teórico del libro de códigos, pero no se someten a evaluación automática.

Esta restricción responde asimismo a una consideración de diseño experimental:
evaluar cinco variables sobre la totalidad del subconjunto anotado y con varios
proveedores de modelos resulta preferible, en términos de solidez de las
conclusiones, a una cobertura amplia pero superficial de las quince.

\subsubsection{Definición operativa de las variables evaluadas}
\label{sec:iris_variables}

Las cinco variables seleccionadas, definidas a partir de las guías de lenguaje no
sexista y del conocimiento experto del proyecto, son las siguientes:

\begin{itemize}
    \item \textbf{V25 --- \texttt{lenguaje\_sexista}.} Uso discriminatorio del
    lenguaje por razón de sexo. Opera como categoría agregadora: engloba todo
    sexismo lingüístico de forma, de modo que la detección de cualquiera de sus
    subtipos debería implicar su marcado. Admite un tercer valor específico para
    el \emph{salto semántico}, esto es, el caso en que un masculino presentado
    como genérico se revela posteriormente no inclusivo al aparecer una marca
    explícita de mujer.

    \item \textbf{V26 --- \texttt{masc\_generico}.} Empleo del masculino
    gramatical como pretendido genérico cuando existe alternativa inclusiva.
    Constituye un subtipo específico de V25.

    \item \textbf{V30 --- \texttt{sexismo\_discurso}.} Sexismo situado en el
    \emph{fondo} del mensaje —en aquello que se cuenta— y no en la forma
    lingüística elegida. Es la única de las cinco que no constituye un subtipo
    de V25, lo que permite que una pieza presente una variable sin la otra.

    \item \textbf{V33 --- \texttt{asimetria\_mujer\_hombre}.} Tratamiento
    asimétrico de mujeres y hombres dentro de una misma pieza: diferencias en el
    uso del nombre, el apellido, el título profesional o el tratamiento.

    \item \textbf{V35 --- \texttt{denominacion\_sexualizada}.} Identificación de
    la mujer por su aspecto físico, su rol familiar o su relación con un varón,
    en detrimento de su función profesional o del motivo informativo por el que
    aparece en la pieza.
\end{itemize}

Todas las variables se codifican en una escala en la que el valor~1 indica
ausencia del fenómeno y el valor~2 su presencia; V25 admite adicionalmente el
valor~3 para el caso del salto semántico. En el tratamiento de los resultados,
y con objeto de homogeneizar la comparación entre variables, los valores 2 y 3
se agrupan como clase positiva.

\subsubsection{Prevalencia y sus implicaciones}

La distribución de las clases en la anotación experta resulta marcadamente
desigual y condiciona de forma decisiva la elección de métricas de evaluación,
tal como se analiza en la Sección~\ref{subsec:iris_divergencia}. La
Tabla~\ref{tab:iris-prevalencia-imp} recoge la prevalencia de la clase positiva de
cada variable y, junto a ella, la exactitud que alcanzaría un clasificador trivial
que respondiese siempre la clase mayoritaria.

\begin{table}[H]
    \centering
    \begingroup
    \small
    \ttabbox[\FBwidth]{
        \caption{Prevalencia de la clase positiva y exactitud del clasificador trivial ($N=1.313$)}
        \label{tab:iris-prevalencia-imp}
    }{
        \begin{tabularx}{\textwidth}{
            l
            X
            >{\centering\arraybackslash}X
            >{\columncolor{gray!15}\centering\arraybackslash}X
        }
            \hline
            \textbf{Cód.} & \textbf{Variable} & \textbf{Prevalencia «Sí»} & \textbf{Exactitud trivial} \\
            \hline
            V25 & \texttt{lenguaje\_sexista}         & 0,810 & 0,810 \\
            V26 & \texttt{masc\_generico}            & 0,810 & 0,810 \\
            V30 & \texttt{sexismo\_discurso}         & 0,428 & 0,572 \\
            V33 & \texttt{asimetria\_mujer\_hombre}  & 0,062 & 0,938 \\
            V35 & \texttt{denominacion\_sexualizada} & 0,100 & 0,900 \\
            \hline
        \end{tabularx}
    }
    \endgroup
\end{table}

La lectura de la última columna es reveladora: en V33 y V35, un modelo que no
detecte nada y responda siempre «No» acertaría el 93,8\,\% y el 90,0\,\% de las
veces. Por consiguiente, \textbf{la exactitud no constituye por sí sola una
métrica informativa} en esta tarea. Todos los resultados de este trabajo se
acompañan del coeficiente $\kappa$ de Cohen~\cite{cohen1960}, que descuenta el acuerdo esperable por
azar, de métricas F1 y de la prevalencia de cada variable.

\subsubsection{De la definición experta a la especificación ejecutable}

Para su empleo por sistemas automáticos, cada variable, partiendo de las guías de
lenguaje no sexista y del criterio experto, se formaliza en una especificación
estructurada que recoge, además de la definición conceptual: un criterio operativo
de decisión, una metodología paso a paso, ejemplos positivos y contraejemplos
comentados, casos límite documentados y las fronteras que la separan de las
variables contiguas. Esta especificación constituye la fuente
única a partir de la cual se generan tanto los \textit{prompts} del sistema de
referencia como las \textit{skills} de la arquitectura de agentes, garantizando
que ambos niveles experimentales operan con \emph{idéntico} contenido
metodológico y difieren únicamente en la forma en que dicho contenido se pone a
disposición del modelo. Los detalles de esta formalización se abordan en el
Capítulo~\ref{implementation}.



\section{Solución tecnológica basada en Inteligencia Artificial}
\label{sec:ai_solution}
%----------
%	SOLUCIÓN TECNOLÓGICA BASADA EN IA (intro de 3.2)
%----------

La tarea descrita en la Sección~\ref{sec:experts} pide clasificar cada pieza
periodística según cinco variables de sexismo lingüístico. Dos rasgos de esa
tarea determinan la aproximación tecnológica. El primero: no hay un conjunto de
entrenamiento etiquetado grande, porque la anotación experta es cara y cubre
1.313 piezas. El segundo: codificar no consiste en buscar palabras clave, sino
en aplicar una metodología interpretativa (definiciones, reglas de decisión,
fronteras entre variables) tomada del libro de códigos del proyecto.

Estos dos rasgos desaconsejan el aprendizaje supervisado clásico, que necesitaría
miles de ejemplos etiquetados. En su lugar se recurre a modelos generativos de
lenguaje (LLM) en régimen \textit{zero-shot}: el modelo no aprende la tarea a
partir de ejemplos, sino que recibe la metodología en el propio \textit{prompt} y
la aplica sobre cada texto. El libro de códigos hace de supervisión en tiempo de
inferencia, sin reentrenar nada.

La solución es una tubería común, independiente del proveedor. Por cada pieza y
variable, el sistema construye una entrada con la metodología y el texto, la
envía a un modelo y recupera una salida estructurada con tres campos: un código,
una explicación razonada y las evidencias literales que la sustentan. Esa misma
tubería se enruta, según el identificador del modelo, tanto a modelos locales
como a las API de OpenAI, Google y Anthropic (Sección~\ref{sec:llms_models_iris}).
Así, las diferencias que se observen entre modelos no se confunden con
diferencias de implementación.

\begin{figure}[H]
\centering
\begin{tikzpicture}[
  node distance=8mm and 12mm,
  caja/.style={draw, rounded corners=3pt, align=center, minimum height=12mm,
               inner sep=5pt, font=\small},
  io/.style={caja, fill=black!5},
  proc/.style={caja, fill=black!2},
  arr/.style={-{Latex[length=2.4mm]}, thick, gray}
]
  \node[io] (art) {Artículo};
  \node[proc, right=of art] (prompt) {Metodología\\ + texto};
  \node[proc, right=of prompt] (modelo) {Modelo\\ (LLM)};
  \node[io, right=of modelo, align=left] (out)
        {\texttt{codigo}\\ \texttt{explicacion}\\ \texttt{evidencias}};
  \draw[arr] (art) -- (prompt);
  \draw[arr] (prompt) -- (modelo);
  \draw[arr] (modelo) -- (out);
\end{tikzpicture}
\caption{Tubería de clasificación. Por cada pieza y variable, la metodología del
libro de códigos y el texto forman la entrada del modelo, que devuelve una salida
estructurada y auditable.}
\label{fig:iris-pipeline}
\end{figure}

La salida estructurada responde al marco \textit{human-in-the-loop} del proyecto.
Pedir explicación y evidencias no es un capricho de formato: convierte cada
predicción en algo que la analista puede auditar, aceptar o rechazar. El sistema
asiste; no decide en su lugar.

Las secciones siguientes describen los modelos generativos considerados
(\ref{sec:llms_models_iris} y \ref{sec:genai_models_iris}) y el paradigma
\textit{zero-shot} con libro de códigos (\ref{sec:zeroshot_iris}). Las dos
arquitecturas de inferencia que se comparan, inyectar la metodología en el
\textit{prompt} frente a empaquetarla como \textit{Agent Skills}, se detallan en
el Capítulo~\ref{implementation}.


\subsection{Modelos generativos de lenguaje}
\label{sec:llms_models}
%----------
%	MODELOS GENERATIVOS DE LENGUAJE
%----------

\label{sec:llms_models_iris}

Un modelo generativo de lenguaje, o LLM, es una red
neuronal, habitualmente de tipo \textit{Transformer}~\cite{vaswani2017}, entrenada
sobre grandes volúmenes de texto para predecir la continuación de una secuencia. De ese entrenamiento sale una competencia
lingüística general que puede orientarse a tareas concretas mediante
instrucciones en lenguaje natural, sin entrenar un modelo distinto para cada
tarea. Este trabajo explota justamente esa competencia: la metodología de
codificación se expresa como instrucción y el modelo la aplica a cada pieza.

Para desplegar estos modelos hay dos vías, y la diferencia entre ellas importa en
un proyecto como este.

Los \textbf{modelos locales} son modelos de pesos abiertos que se ejecutan en
infraestructura propia. En este trabajo se usan a través de
\textit{Ollama}\footnote{\url{https://ollama.com} [Último acceso: 05/08/2026]} sobre la granja de GPU del
Departamento de Teoría de la Señal y Comunicaciones. Tienen
dos ventajas principales. La primera es la privacidad: el texto nunca sale de la
infraestructura de la institución, algo relevante cuando se trabaja con contenido
periodístico y datos de investigación. La segunda es el coste: la inferencia no se
paga por uso, a diferencia de las API comerciales, por lo que procesar el corpus
completo no tiene coste marginal. A cambio, exigen hardware especializado y suelen
ser más pequeños y menos capaces que los modelos comerciales punteros.

Los \textbf{modelos comerciales vía API} los alojan las empresas proveedoras (OpenAI,
Google) y se acceden por red. Ofrecen más capacidad y no requieren
infraestructura propia. La contrapartida es doble: se paga por uso y el texto
viaja a un tercero, con lo que ello implica en privacidad.

Comparar ambas vías sobre la misma tarea y el mismo corpus es uno de los
objetivos metodológicos de este trabajo. Para lograrlo, la cadena de procesos de
clasificación es agnóstica a la proveedora: una función de enrutado dirige cada
petición al \textit{backend} que corresponda según el identificador del modelo. De este
modo, lo que se observe al comparar modelos procede de los modelos, no de cómo
está escrito el código que los llama. Los modelos concretos se detallan en la
Sección~\ref{sec:genai_models_iris}.


\subsection{Modelos evaluados}
\label{genai_models}
%----------
%	MODELOS EVALUADOS (multi-proveedor)
%----------

\label{sec:genai_models_iris}

La selección cubre las dos vías de despliegue y las principales familias
comerciales. En cada caso se eligen variantes ligeras y económicas, aptas para
procesar un corpus periodístico completo. La Tabla~\ref{tab:iris-modelos} recoge
los modelos considerados.

\begin{table}[H]
    \centering
    \begingroup
    \small
    \ttabbox[\FBwidth]{
        \caption{Modelos considerados y vía de acceso}
        \label{tab:iris-modelos}
    }{
        \begin{tabularx}{\textwidth}{
            l
            l
            >{\raggedright\arraybackslash}X
        }
            \hline
            \textbf{Proveedor} & \textbf{Modelo} & \textbf{Vía / observaciones} \\
            \hline
            Local (Ollama) & \texttt{gemma4:e4b} & Granja GPU TSC; modelo abierto de Google \\
            OpenAI & \texttt{gpt-4o-mini} & API \\
            OpenAI & \texttt{gpt-5.4-nano} & API; variante de razonamiento (desactivado) \\
            Google & \texttt{gemini-3.1-flash-lite} & API \\
            \hline
        \end{tabularx}
    }
    \endgroup
\end{table}

La evaluación del Capítulo~\ref{results_and_discussion} cubre estos cuatro
modelos: \texttt{gemma4:e4b} en local y las tres variantes comerciales de API.
Tres decisiones sobre la selección merecen justificarse.

Para que la comparación entre proveedores sea justa, se desactivó el razonamiento
explícito en los modelos que lo admiten (el parámetro \texttt{reasoning\_effort}
en OpenAI, \texttt{thinking\_budget} en Google). Así ningún modelo dispone de
cómputo extra del que carezcan los demás.

Los modelos locales no cuestan dinero, porque corren en infraestructura propia,
pero sí tienen un coste computacional alto que condiciona su viabilidad a escala
(Sección~\ref{subsec:iris_coste}).

La preferencia por variantes ligeras (\textit{mini}, \textit{nano},
\textit{flash-lite}, modelos locales de pocos miles de millones de parámetros) es
deliberada. Un sistema de asistencia editorial que procese miles de piezas
necesita un coste por documento contenido. Esa elección tiene un precio, menor
capacidad frente a los modelos punteros, que se discute a la luz de los
resultados en el Capítulo~\ref{results_and_discussion}.


\subsection{Clasificación zero-shot con libro de códigos}
\label{zero_shot}
%----------
%	CLASIFICACIÓN ZERO-SHOT CON LIBRO DE CÓDIGOS
%----------

\label{sec:zeroshot_iris}

La clasificación se aborda en régimen \textit{zero-shot}, presentado en la
Sección~\ref{sec:ai_solution}. Al no disponerse de un corpus de entrenamiento
etiquetado a gran escala, la supervisión por ejemplos se sustituye por la
\textbf{supervisión mediante libro de códigos}~\cite{ziems2024css, hollauer_2024},
de modo que el modelo generaliza a partir del conocimiento que ya posee y de la
metodología que recibe en la instrucción.

La metodología que se traslada al \textit{prompt} procede de la especificación
estructurada de cada variable descrita en la Sección~\ref{sec:iris_variables}:
criterio de decisión, pasos, ejemplos y fronteras. Todo ese conocimiento experto
se centraliza en un único fichero, \texttt{variables.json},
que actúa como fuente de verdad del sistema. A partir de él se generan de forma
automática las dos arquitecturas que se comparan en el
Capítulo~\ref{implementation}, de modo que ambas parten de idéntico contenido y
difieren solo en cómo lo reciben.

\subsubsection{Salida estructurada y trazabilidad}

Por cada pieza y variable, el modelo devuelve un JSON con tres campos:

\begin{itemize}
    \item \texttt{codigo}: el valor en la escala cruda de la variable
    (1 = ausencia, 2 = presencia, y V25 admite además el 3 para el salto semántico).
    El modelo emite este código sin agrupar. Solo al calcular las métricas los
    valores 2 y 3 se unen como clase positiva (Sección~\ref{sec:iris_variables}).
    \item \texttt{explicacion}: la justificación, que debe citar los pasos de la
    metodología aplicados.
    \item \texttt{evidencias}: fragmentos literales del texto que sostienen la
    decisión (lista vacía cuando \texttt{codigo}\,=\,1, es decir, en los negativos).
\end{itemize}

\begin{figure}[H]
    \centering
    \includegraphics[width=\textwidth]{Plantilla_TFG_ingles_2019/imagenes/salida_trazabilidad_blanco.pdf}
    \caption{Salida estructurada y trazable de una predicción.}
    \label{fig:iris-salida}
\end{figure}

Este formato materializa el marco
\textit{human-in-the-loop}. La explicación y las evidencias hacen la predicción
auditable: la analista lee la justificación y decide si acepta, rechaza o matiza
el veredicto. Detectar un sesgo de género no se reduce a localizar palabras
clave, sino que exige interpretar el contexto social del texto, un criterio que
sigue siendo competencia humana. Por eso el sistema asiste a la analista en lugar
de sustituirla.

\subsubsection{Ventajas y límites del enfoque}

El alcance de esta decisión se aprecia mejor si se examinan sus efectos sobre dos
planos distintos, el de la adquisición de conocimiento y el del modelado
estadístico de la tarea.

En el plano de la \textbf{adquisición de conocimiento}, el sistema no incorpora
nada de lo que procesa. Los pesos del modelo permanecen intactos, de modo que la
metodología no llega a formar parte de él, sino que se le suministra en cada
llamada y se pierde al terminarla. Esto tiene una cara favorable y otra adversa.
La favorable es que prescinde de conjuntos de entrenamiento etiquetados y del
preprocesamiento intensivo de los enfoques clásicos, y sobre todo que mejorar el
sistema consiste en reescribir la metodología y no en reentrenar un modelo, tarea
que corresponde entonces al equipo experto y no a quien programa. La adversa es
que el sistema no acumula experiencia, ya que un error cometido en una pieza no
informa el análisis de la siguiente, y todo el conocimiento del dominio con que
cuenta procede de su preentrenamiento, ajeno a este corpus y a esta metodología.

En el plano del \textbf{modelado estadístico}, el sistema no estima probabilidad
alguna a partir de los datos anotados. Un clasificador supervisado aprendería con
qué frecuencia aparece cada clase y ajustaría en consecuencia su umbral de
decisión, de manera que en una variable como V33, con un 6,2\,\% de positivos,
tendería por sí solo a marcar poco. Aquí no sucede nada de eso. Al no observar
nunca la distribución de clases del corpus (Sección~\ref{sec:iris_variables}), la
propensión del modelo a marcar no responde a la prevalencia real, sino a cómo
interpreta el criterio que se le ha descrito y a las regularidades de su
preentrenamiento. Cuando marca de más o de menos no está, por tanto, mal
calibrado respecto a una distribución que conoce, sino aplicando un umbral
interpretativo propio, observación que resulta determinante para leer los
resultados del Capítulo~\ref{results_and_discussion} y, en particular, las
diferencias de recall entre modelos.

De ambos planos se sigue la limitación principal del enfoque. El rendimiento queda
determinado por la forma en que la metodología se traslada al \textit{prompt} y
por la capacidad del modelo empleado, sin que existan ejemplos etiquetados con los
que corregir sus desviaciones. Determinar el alcance de esa limitación, y
establecer si su origen reside en el modelo, en la forma de presentarle la
metodología o en la definición misma de la tarea, constituye el objeto del
Capítulo~\ref{results_and_discussion}.



    %----------
    %	IMPLEMENTATION
    %----------	

\chapter{Implementación}
\label{implementation}
Este capítulo describe cómo se lleva a la práctica la solución planteada en el
Capítulo~\ref{sec:methodology}. Se detallan la arquitectura de agentes y sus dos
niveles de comparación (Sección~\ref{agents}), las \textit{skills} y herramientas
que la sostienen (Sección~\ref{prompts}), la recuperación sobre las guías
expertas (Sección~\ref{rag}), la instrumentación de la evaluación
(Sección~\ref{eval}) y la prueba de concepto de aplicabilidad editorial
(Sección~\ref{poc}).

La infraestructura combina dos entornos, unidos por una tubería común y agnóstica
al proveedor. Los modelos comerciales se acceden por API (OpenAI, Google). Los
modelos de pesos abiertos se ejecutan mediante \textit{Ollama}
sobre la granja de GPU del Departamento de Teoría de la Señal y Comunicaciones, a
la que se envían las peticiones desde un equipo cliente del laboratorio. El código
fuente completo se recoge en el Apéndice~\ref{appendix_code}.







\section{Arquitectura y diseño agéntico}
\label{agents}
%----------
%	ARQUITECTURA Y DISEÑO AGÉNTICO
%----------

El núcleo de esta investigación compara dos formas de poner la metodología experta a
disposición del modelo. Ambas resuelven la misma tarea, sobre el mismo corpus y
con el mismo número de llamadas al modelo. Lo único que cambia entre ellas es
cómo llega esa metodología.

\begin{itemize}
    \item \textbf{Nivel B0 (\textit{baseline}).} La metodología de cada variable
    se inyecta entera en el \textit{prompt}. El modelo emite su veredicto en una
    sola llamada, sin herramientas.
    \item \textbf{Nivel B1 (\textit{Agent Skills}).} El modelo recibe únicamente
    los metadatos de las \textit{skills} disponibles (su nombre y una
    descripción) y carga bajo demanda el cuerpo de la que necesita. Este mecanismo
    se conoce como \textit{progressive disclosure}. Además puede consultar las
    guías expertas y verificar sus evidencias mediante herramientas.
\end{itemize}

\begin{figure}[H]
\centering
\includegraphics[width=0.94\textwidth]{arquitectura_b0_b1_blanco.pdf}
\caption{Los dos niveles comparados. Ambos hacen cinco llamadas por artículo (una
por variable) y comparten formato de salida.}
\label{fig:iris-b0b1}
\end{figure}

El diseño fija todo salvo un factor: ambos niveles hacen \textbf{cinco llamadas
por artículo}, una por variable, y comparten el mismo formato de salida. Lo único
que cambia es cómo llega la metodología al modelo. De este modo, el delta de
acuerdo entre B0 y B1 es atribuible al empaquetado como \textit{Agent Skills}, y no
al número de llamadas ni al formato de la respuesta.

\subsection{Diseño agéntico: un agente por dentro}
\label{subsec:diseno_agentico}

Antes de describir cómo se organiza la arquitectura  multi-agéntica, conviene fijar de qué se
compone \emph{uno}. Para ello hay que distinguir dos conceptos que a menudo se
confunden. Una \textit{skill}\footnote{El concepto de \textit{Agent Skills} fue
formalizado por Anthropic a finales de 2025 como un estándar abierto para la
portabilidad entre plataformas: \url{https://www.anthropic.com/engineering/equipping-agents-for-the-real-world-with-agent-skills} [Último acceso: 05/08/2026]}
es conocimiento empaquetado: un fichero de metodología, pasivo, que se carga pero
no decide. Un \textbf{agente} es un ejecutor con un bucle de decisión, que elige
qué skills cargar y cuándo verificar. \textit{Una skill por variable no es un agente} pero
\textit{un agente que carga su skill bajo demanda} sí lo es, y es lo que aquí se
implementa.

Cada agente dispone de un \textbf{catálogo de skills} del que, de entrada, solo ve
el nombre y una descripción; carga el cuerpo de una skill únicamente cuando la
necesita. Este mecanismo se conoce como \textit{progressive disclosure}, un
principio del diseño de interacción persona-máquina (HCI por sus siglas en inglés), mostrar solo lo esencial y revelar el detalle
bajo demanda, que se popularizó en el ámbito de la
usabilidad~\cite{nielsen1994} y que la comunidad de agentes ha tomado
prestado para el diseño del contexto. La idea es la misma: cargar únicamente lo
relevante para la tarea en curso evita saturar la ventana de contexto con
información que no viene al caso, algo que degrada la respuesta y encarece la
inferencia. No procede de una fuente académica primaria sobre agentes, sino de esa
tradición de usabilidad; su formulación operativa para modelos es la del marco de
\textit{Agent Skills} citado más arriba. El catálogo reúne tres tipos de skill: la
de la propia variable (obligatoria), unas auxiliares compartidas y unos resúmenes
de las guías expertas, cuyo formato y generación se detallan en la
Sección~\ref{prompts}.

El agente opera mediante un \textbf{bucle estilo ReAct}~\cite{yao2023react} (de
\textit{Reasoning} y \textit{Acting}), un patrón que entrelaza razonamiento y
acción en vez de pedir la respuesta de un solo tirón. En cada turno el modelo razona sobre qué le falta,
emite \emph{una} acción, observa el resultado que le devuelve el sistema y vuelve a
razonar con esa nueva información; el ciclo se repite hasta que cierra con su
veredicto. Ese entrelazado es lo que le permite, por ejemplo, cargar su skill,
consultar una guía si duda y verificar sus evidencias antes de responder, en lugar
de contestar a ciegas.

Las acciones se expresan \textbf{en texto}, no mediante el \textit{tool-calling}
nativo de cada proveedor. La razón es práctica: el enrutado multi-proveedor se
apoya en una interfaz de texto común, y el \textit{tool-calling} nativo es frágil e
incompatible entre modelos, sobre todo en los locales pequeños. Con un protocolo de
texto, el mismo mecanismo funciona igual en los cuatro proveedores. En cada turno
el agente emite una de estas acciones y el sistema le devuelve el resultado:


\begin{itemize}
    \item \texttt{LEER\_SKILL <id>}: carga el cuerpo de una skill.
    \item \texttt{CONSULTAR\_GUIA <consulta>}: recupera pasajes literales de las
    guías expertas (Sección~\ref{rag}).
    \item \texttt{VERIFICAR ["..."]}: comprueba que las evidencias son literales
    del texto.
    \item \texttt{FINAL \{...\}}: cierra con el veredicto de la variable.
\end{itemize}

El bucle admite hasta unas ocho iteraciones. Si el agente emitiera \texttt{FINAL}
sin haber usado ninguna herramienta, su comportamiento equivaldría al del
\textit{baseline}; ese «colapso a B0» se registra como una métrica más, aunque en
la práctica apenas se observa, porque todos los agentes cargan al menos su propia
skill. Más informativa resulta la \emph{tasa de uso} de cada herramienta: con qué
frecuencia un modelo consulta las guías o verifica sus evidencias. Esa medida de
comportamiento se reporta en el Capítulo~\ref{results_and_discussion} y ayuda a
interpretar por qué la arquitectura mejora el acuerdo en unos modelos y no en
otros.

\begin{figure}[H]
\centering
\includegraphics[width=\textwidth]{agente_anatomia_blanco.pdf}
\caption{Anatomía de un agente especializado, con el ejemplo del agente de
\texttt{lenguaje\_sexista}: catálogo de skills (progressive disclosure) y bucle de
acciones que carga la skill bajo demanda, consulta guías, verifica evidencias y
cierra con el veredicto.}
\label{fig:iris-agente-anatomia}
\end{figure}

\subsection{Arquitectura multiagente}
\label{subsec:multiagente}

Sobre ese diseño se construye la arquitectura del nivel B1: no un único agente que
juzgue las cinco variables a la vez, sino \textbf{cinco agentes especializados},
uno por variable (V25, V26, V30, V33 y V35). Cada uno tiene un único objetivo,
clasificar su variable, y resuelve su tarea de forma independiente de los demás,
con su propio catálogo y su propio bucle. Cada artículo se procesa, por tanto, con
cinco agentes y cinco llamadas, una por variable.

La especialización responde a una decisión de diseño. Sobrecargar a un modelo con
mucho contexto tiende a degradar su rendimiento~\cite{liu2023lost}: cuando la
información relevante se diluye entre la que no lo es, aumentan las
alucinaciones~\cite{huang2025hallucination} y la confusión entre criterios. Un solo agente encargado de las cinco variables tendría que sostener a
la vez cinco metodologías y sus fronteras, incurriendo en ese riesgo. Acotar cada
agente a una variable reduce la carga de contexto, mantiene el foco y hace que su
traza (qué skills cargó, si verificó, si colapsó a B0) sea interpretable variable a
variable.

\begin{figure}[H]
\centering
\includegraphics[width=\textwidth]{agentes_enjambre_blanco.pdf}
\caption{Arquitectura multiagente del nivel B1: cada artículo se procesa con cinco
agentes especializados independientes, uno por variable, cada uno con su propio
catálogo y su propio bucle.}
\label{fig:iris-multiagente}
\end{figure}


\section{Prompts y herramientas de análisis}
\label{prompts}
%----------
%	SKILLS, PROMPTS Y HERRAMIENTAS
%----------

Toda la metodología del sistema vive en ficheros \texttt{SKILL.md}: un documento
por unidad de conocimiento, con un encabezado (nombre y descripción) y el cuerpo
con las instrucciones. La descripción es lo que el agente ve antes de decidir si
carga la skill, y por tanto lo que gobierna el enrutado. El cuerpo solo se lee
cuando hace falta.

\subsection{Tipos de skill}

Hay tres tipos, con papeles distintos:

\begin{itemize}
    \item \textbf{Skills de variable} (cinco). Contienen la metodología paso a
    paso de cada variable. Se generan automáticamente desde
    \texttt{variables.json}, de modo que el codebook y las skills nunca se
    desincronizan. Cada agente carga la suya.
    \item \textbf{Skills auxiliares} (tres): \texttt{guia\_regla\_inversion},
    \texttt{guia\_lenguaje\_inclusivo} y \texttt{verificar\_evidencias}.
    Cualquier agente las carga bajo demanda cuando duda.
    \item \textbf{Resúmenes de guías}. Condensaciones de las guías reales de
    lenguaje no sexista (Sección~\ref{rag}), disponibles como orientación rápida.
\end{itemize}

Que las skills de variable se generen desde \texttt{variables.json} tiene una
consecuencia importante para la validez del experimento: el nivel B0 y el nivel
B1 parten del mismo fichero fuente. El \textit{baseline} recibe ese contenido
inyectado en el \textit{prompt}; el agente lo carga como skill. El contenido es
byte a byte el mismo, así que la comparación aísla el mecanismo de entrega y no
un cambio en la metodología.

\subsection{Herramientas}

Las herramientas del agente son deliberadamente pocas y de propósito claro:

\begin{itemize}
    \item \texttt{list\_skills} devuelve solo los metadatos (nombre y
    descripción) de las skills visibles. Es lo que hace efectiva la
    \textit{progressive disclosure}: el agente elige a ciegas del cuerpo.
    \item \texttt{read\_skill} entrega el cuerpo completo de una skill concreta.
    \item \texttt{verificar\_evidencias} comprueba, contra el mismo texto que vio
    el modelo, qué fragmentos citados son literales.
\end{itemize}

\subsection{Salida estructurada y verificación de evidencias}

Cada agente cierra con un objeto JSON de tres campos (\texttt{codigo},
\texttt{explicacion}, \texttt{evidencias}), descrito en la
Sección~\ref{sec:zeroshot_iris}. Sobre las evidencias se aplica una comprobación
automática: solo se conservan los fragmentos que aparecen \emph{literalmente} en
el texto original. Esta verificación cumple dos funciones. Refuerza la
trazabilidad exigida por el marco \textit{human-in-the-loop}, y evita un fallo
sutil que arrastraban implementaciones anteriores, donde el texto se alteraba
antes de contrastar las evidencias y estas se descartaban por error. Aquí la
comprobación se hace contra el texto exacto que recibió el modelo.


\section{Representaciones: RAG, recuperación y similitud}
\label{rag}
%----------
%	RECUPERACIÓN SOBRE GUÍAS EXPERTAS (RAG)
%----------

El codebook del proyecto se inspira en guías institucionales de lenguaje no
sexista (Sección~\ref{sec:experts}). Esas guías cumplen dos papeles distintos en
el sistema, y conviene no confundirlos.

Como \textbf{fuente}, son el material del que se destila la metodología de cada
variable. Ese uso es indirecto: la analista y quien redacta las skills se apoyan
en ellas, pero el modelo no las lee.

Como \textbf{herramienta en vivo}, en cambio, el agente puede consultarlas
durante la inferencia. La acción \texttt{CONSULTAR\_GUIA} implementa una
recuperación (\textit{Retrieval-Augmented Generation}, RAG por sus siglas en inglés)~\cite{lewis2020rag} sobre el texto real
de las guías. La implementación es deliberadamente ligera: un índice TF-IDF en
memoria sobre los pasajes de las guías, sin dependencias externas. Ante una
consulta, devuelve los pasajes más relevantes con su cita de fuente y sección.

Esta distinción importa para la honestidad del trabajo. Sin la herramienta, la
conexión con las guías sería solo conceptual, y afirmar que \textit{``el agente se apoya en las guías expertas''} sería inexacto. Con \texttt{CONSULTAR\_GUIA}, esa afirmación
se sostiene: el agente recupera texto literal de las fuentes en tiempo de
inferencia y puede citarlo en su justificación.

La recuperación se basa en coincidencia léxica (TF-IDF)~\cite{sparckjones1972}, no semántica. Funciona
bien para términos precisos, como \textit{masculino genérico} o \textit{feminización de
cargos}, pero es sensible a sinónimos y paráfrasis. Una recuperación por
\textit{embeddings} mejoraría esa robustez y queda como línea de mejora. Para el
alcance de este trabajo, el índice léxico es suficiente. Cuestión distinta es si esa recuperación mejora el acuerdo final: su
efecto se evalúa por separado en la Sección~\ref{subsec:iris_descomposicion}, y
resulta, en la práctica, prescindible.


\section{Evaluación de modelos resultantes}
\label{eval}
%----------
%	EVALUACIÓN E INSTRUMENTACIÓN
%----------

\subsection{Métricas}

Las predicciones se comparan con la anotación experta sobre las 1.313 piezas
etiquetadas. Cada variable se trata como un problema binario en el que la clase
positiva es la presencia del fenómeno, es decir, el \textit{Sí}. De esa
comparación resultan cuatro recuentos: los verdaderos positivos ($VP$), piezas que
tanto el sistema como la anotadora marcan, los falsos positivos ($FP$), que el
sistema marca y la anotadora no, los falsos negativos ($FN$), que la anotadora
marca y el sistema pasa por alto, y los verdaderos negativos ($VN$), en que ambas
coinciden en la ausencia. Sobre esos cuatro valores se construyen todas las
métricas empleadas.

\paragraph{Exactitud.} Es la proporción de piezas en que sistema y anotadora
coinciden, con independencia de la clase:
\begin{equation}
    \text{Exactitud} = \frac{VP + VN}{VP + VN + FP + FN}
\end{equation}
Su ventaja es que resulta inmediatamente interpretable, ya que responde a la
pregunta de cuántas veces acierta el sistema. Su inconveniente es que la domina la
clase mayoritaria. En V33, con un 6,2\,\% de positivos, no marcar nada produce una
exactitud del 93,8\,\% sin haber detectado un solo caso
(Sección~\ref{sec:iris_variables}), de manera que por sí sola no distingue un
sistema competente de uno que no lo es.

\paragraph{Precisión.} De todo lo que el sistema marca, mide qué proporción estaba
efectivamente marcada por la anotadora:
\begin{equation}
    \text{Precisión} = \frac{VP}{VP + FP}
\end{equation}
Responde a la pregunta de cuánto cabe fiarse de una alerta del sistema, lo que la
hace directamente relevante para una herramienta de asistencia, ya que una
precisión baja se traduce en avisos que la persona analista deberá descartar. Su
inconveniente es que se calcula solo sobre lo marcado, así que un sistema muy
conservador puede alcanzar una precisión alta con un puñado de aciertos, y
queda indefinida si no marca nada.

\paragraph{Recall.} También llamado sensibilidad, mide qué proporción de lo que la
anotadora marcó consigue detectar el sistema:
\begin{equation}
    \text{Recall} = \frac{VP}{VP + FN}
\end{equation}
Responde a la pregunta complementaria, la de cuánto sexismo se está dejando sin
señalar. Su inconveniente es simétrico al de la precisión, puesto que marcarlo
todo garantiza un recall de 1 sin discriminar nada. Precisión y recall solo tienen
sentido, por tanto, leídos conjuntamente.

\paragraph{F1.} Resume ambas en su media armónica, que penaliza el desequilibrio
entre las dos más de lo que lo haría una media aritmética:
\begin{equation}
    F_1 = 2 \cdot \frac{\text{Precisión} \cdot \text{Recall}}{\text{Precisión} + \text{Recall}}
\end{equation}
La variante \textbf{macro} calcula el $F_1$ de cada clase por separado y promedia
los dos valores sin ponderar por su frecuencia, de modo que la clase minoritaria
pesa tanto como la mayoritaria. Es una corrección importante frente a la
exactitud, pero no elimina el problema del desequilibrio, ya que en una variable
con un 81\,\% de positivos como V25 marcarlo todo produce un $F_1$ de 0,895 sobre
la clase positiva.

\paragraph{Coeficiente $\kappa$ de Cohen.} Mide el acuerdo entre dos anotadores
descontando el que cabría esperar por azar~\cite{cohen1960}:
\begin{equation}
    \kappa = \frac{p_o - p_e}{1 - p_e}
\end{equation}
donde $p_o$ es el acuerdo observado, equivalente a la exactitud, y $p_e$ el
acuerdo esperado si ambos etiquetasen al azar respetando sus respectivas
frecuencias marginales. Un valor de 1 indica acuerdo perfecto, un 0 indica un
acuerdo equivalente al del azar y los valores negativos señalan un desacuerdo
sistemático. Su ventaja es que sitúa el resultado frente a una referencia
razonable en lugar de frente al cero absoluto, y que procede del análisis de
contenido, disciplina de la que nace este trabajo. Su inconveniente es la conocida
\textbf{paradoja de $\kappa$}~\cite{feinstein1990}, según la cual con
distribuciones muy desequilibradas el acuerdo esperado por azar se dispara y el
coeficiente se desploma aunque la exactitud sea alta. Esto vuelve los valores de
$\kappa$ difícilmente comparables entre variables de prevalencia distinta.

Ninguna métrica basta por sí sola, y la que resulta informativa depende de la
variable. En V33 y V35, donde los positivos son escasos, la exactitud carece de
valor, la precisión se vuelve inestable por apoyarse en muy pocos casos y el
recall es el indicador decisivo, ya que revela cuánto se queda sin detectar. En
V25 y V26, con prevalencias del 81\,\%, ocurre lo contrario, pues el recall y el
$F_1$ se inflan con facilidad y lo informativo pasa a ser la capacidad de
reconocer las piezas sin sexismo, que es donde $\kappa$ resulta más exigente. Por
este motivo los resultados del Capítulo~\ref{results_and_discussion} se presentan
siempre acompañados de la prevalencia de cada variable y del rendimiento del
clasificador trivial (Sección~\ref{sec:iris_variables}), que fijan la referencia
mínima frente a la que interpretarlos.

Existe, además, un criterio que ordena esta elección y que procede del uso
previsto del sistema. Como herramienta de apoyo a la postedición, su cometido
consiste en señalar a la persona analista las piezas que merecen una segunda
lectura, de modo que un falso positivo cuesta unos segundos de descarte mientras
que un falso negativo supone un caso de sexismo que nadie revisará. Esa asimetría
otorga al recall una prioridad que no tendría en un clasificador autónomo, y así se
discute en el Capítulo~\ref{results_and_discussion}.


\subsection{Evaluación \textit{zero-shot} sobre el corpus completo}

La evaluación se realiza en régimen \textit{zero-shot}
(Sección~\ref{sec:zeroshot_iris}), lo que significa que el modelo recibe la
metodología en el \textit{prompt} y que no se ajusta ningún parámetro a la vista
de los datos. Dos consecuencias de ese régimen, ya justificadas en la
Sección~\ref{sec:zeroshot_iris}, condicionan la lectura de los resultados. En
cuanto a la adquisición de conocimiento, el sistema no incorpora información del corpus, de modo que su competencia procede íntegramente del preentrenamiento y de la
metodología que se le entrega en cada llamada. En cuanto al modelado estadístico,
no estima probabilidades a partir de los datos anotados ni ajusta su umbral a la
prevalencia de cada variable, así que su propensión a marcar refleja su
interpretación del criterio y no la distribución real de las clases.

De ahí se desprende una consecuencia práctica sobre el diseño de la evaluación.
Todas las cifras del Capítulo~\ref{results_and_discussion} se calculan sobre las
1.313 piezas etiquetadas, sin reservar un subconjunto de desarrollo o de prueba.
Al no existir entrenamiento ni selección de umbrales sobre los datos, no cabe el
riesgo de sobreajustar a un subconjunto y medir después el rendimiento en él, y el
corpus completo puede emplearse como material de evaluación.


\subsection{Ejecución a escala: enrutado y granja}

La misma secuencia de procesos sirve para todos los modelos gracias al enrutamiento por
identificador (Sección~\ref{sec:llms_models_iris}). Para los modelos comerciales,
las peticiones van a la API correspondiente. Para los locales, se ejecutan sobre
la granja de GPU del Departamento de Teoría de la Señal y Comunicaciones mediante
\textit{Ollama}, repartiendo el corpus en fragmentos (\textit{shards}) que se procesan
en paralelo sobre varios servidores, con las peticiones lanzadas desde un equipo
cliente del laboratorio. El coste computacional de esta vía local se cuantifica en
la Sección~\ref{subsec:iris_coste}. El código se recoge en el
Apéndice~\ref{appendix_code}.


\section{Prueba de concepto de aplicabilidad editorial}
\label{poc}
%----------
%	PRUEBA DE CONCEPTO DE APLICABILIDAD EDITORIAL
%----------

La finalidad última del sistema no es puntuar en una evaluación, sino asistir el
trabajo de análisis sobre el corpus, incluidas las miles de piezas sin anotar.
Para materializar esa aplicabilidad, el sistema de detección se integra en la
\textbf{herramienta web del proyecto IRIS}~\cite{avi2026}\footnote{\url{https://dei.inf.uc3m.es/iris/} [Último acceso: 05/08/2026]},
y la prueba de concepto consiste en usar esa herramienta para acompañar a expertas
y periodistas en el análisis de un corpus de prueba.

La asistencia cubre los dos planos del análisis de contenido experto
(Sección~\ref{sec:experts}). En el plano \textbf{cuantitativo}, la herramienta agrega
las etiquetas que produce el sistema en informes por variable, con tablas de
frecuencias y porcentajes sobre el conjunto seleccionado y un análisis descriptivo
automático (Figura~\ref{fig:iris-web-informe}), lo que permite explorar qué medios,
secciones o periodos concentran más usos marcados. En el plano \textbf{cualitativo},
para cada pieza muestra el marcado de las evidencias sobre el texto y la explicación
del veredicto (Figura~\ref{fig:iris-web-analisis}), de modo que quien analiza puede
leer por qué el sistema ha decidido lo que ha decidido.

El sistema no automatiza la decisión: la produce en un formato revisable (código,
justificación y evidencias literales) y la analista acepta, rechaza o matiza cada
veredicto.

Los resultados del Capítulo~\ref{results_and_discussion} obligan a matizar el
alcance. El acuerdo entre el sistema y las anotadoras es bajo, y buena parte de ese
desacuerdo podría tener que ver con lo amplio y variable del criterio experto más que con
un fallo de detección. Por eso la herramienta no se plantea como un clasificador
autónomo, sino como un
\textbf{apoyo a la postedición}: en lugar de sustituir el criterio editorial, se
inserta en el flujo de trabajo de quien revisa un texto ya redactado. El sistema
señala las piezas candidatas y aporta en cada caso las evidencias y la
justificación del veredicto, de modo que la analista edita ese diagnóstico
(lo acepta, lo rechaza o lo matiza) en segundos y con plena trazabilidad, en vez
de partir de cero.

Las Figuras~\ref{fig:iris-web-informe} y~\ref{fig:iris-web-analisis} ilustran ambas
vistas de la herramienta web, la de informe cuantitativo y la de análisis por pieza.

\begin{figure}[H]
    \centering
    \includegraphics[width=\textwidth]{Plantilla_TFG_ingles_2019/imagenes/iris_3.png}
    \caption{Vista de informe cuantitativo en la herramienta web del proyecto IRIS.
    Para las variables seleccionadas genera tablas de frecuencias y porcentajes sobre
    el conjunto de piezas, acompañadas de un análisis descriptivo automático.}
    \label{fig:iris-web-informe}
\end{figure}

\begin{figure}[H]
    \centering
    \includegraphics[width=\textwidth]{Plantilla_TFG_ingles_2019/imagenes/iris_2.png}
    \caption{Vista de resultado del análisis en la herramienta web del proyecto IRIS.
    Las detecciones se resaltan sobre el texto y se listan por categoría, con los
    controles para que la analista acepte, edite o rechace cada una (postedición con
    la persona en el bucle).}
    \label{fig:iris-web-analisis}
\end{figure}









    %----------
    %	RESULTS AND DISCUSSION
    %----------	

\chapter{Resultados y Discusión}
\label{results_and_discussion}
%----------
%	RESULTADOS Y DISCUSIÓN — introducción del capítulo
%----------

Este capítulo evalúa la arquitectura de agentes especializadas con \textit{Agent
Skills} presentada en el Capítulo~\ref{implementation}, aplicada a la
clasificación de las cinco variables de lenguaje sexista del corpus definido en el Capítulo~\ref{sec:methodology}. La
Sección~\ref{sec:iris_experimentacion} describe el diseño experimental y presenta
los resultados sobre las 1.313 noticias con anotación experta, mientras que la
Sección~\ref{sec:discussion} los discute de forma transversal y extrae las
implicaciones para el proyecto IRIS. La Sección~\ref{poc} cierra el capítulo
examinando, a la luz de esos resultados, de qué modo el sistema puede emplearse
en el trabajo editorial y cómo se ha integrado en la herramienta del proyecto.

Este capítulo adopta una óptica concreta para interpretar los resultados. El interés
no recae únicamente en el acuerdo entre el sistema y la anotación, sino en lo que el
acuerdo, y sobre todo el desacuerdo, revela sobre la naturaleza de la tarea. La razón
es que el sexismo lingüístico resulta difícil de percibir, y lo resulta también para
quien está formada para percibirlo, porque la ceguera que produce la cultura patriarcal y el espejismo de
neutralidad androcéntrica actúan sobre usos que la lengua ha normalizado hasta
volverlos poco visibles. De ahí que la variabilidad que se observa entre personas
anotadoras no deba leerse sin más como una divergencia de pareceres sobre qué es
sexismo, sino como el reflejo de esa dificultad. Y de ahí también que un acuerdo bajo
admita dos lecturas que el capítulo se ocupa de separar: la de una tarea cuyo criterio
no llega a estabilizarse y la de un sistema que no detecta parte importante de lo que las expertas sí
señalan de forma mayoritaria, con independencia de quién anote. Es importante señalar que cada una de las piezas con las que trabajamos sólo ha sido marcada por una experta sin que exista, por tanto, la posibilidad de verificar el acuerdo entre ellas. 



\section{Detección de lenguaje sexista con agentes especializados}
\label{sec:iris_experimentacion}
% =====================================================================
% Experimentación y resultados — Agentes especializados y Agent Skills
% Datos: Experimentos/experiments/experimento_21_agentskills/results/
% Corpus de evaluación: 1.313 noticias anotadas en las cinco variables
% =====================================================================

La pregunta que vertebra toda la experimentación es una: \textbf{¿aporta algo
empaquetar la metodología experta como \textit{Agent Skills}?} Para responderla se
enfrenta la arquitectura propuesta (\textbf{nivel B1}) a una condición de control
que recibe la misma metodología por la vía más simple, inyectada íntegramente en el
\textit{prompt} (\textbf{nivel B0}). Ambos niveles se definen en el
Capítulo~\ref{implementation} (Sección~\ref{subsec:diseno_agentico}).

Aislar el efecto de las \textit{skills} exige que todo lo demás permanezca fijo.
Por eso B0 y B1 comparten el mismo corpus, los mismos
modelos, las mismas cinco variables y el mismo formato de salida: la única
diferencia entre ambos es cómo llega la metodología al modelo, inyectada de una vez
en el \textit{prompt} (B0) o cargada bajo demanda por el agente (B1), de modo que
cualquier variación en el acuerdo es atribuible al empaquetado como \textit{Agent
Skills} y no a otro factor. Para homogeneizar además el comportamiento entre
proveedores, se desactivó el razonamiento explícito en todos los modelos
(\texttt{reasoning\_effort} en OpenAI, \texttt{thinking\_budget} en Google). La
interpretación de las métricas exige tener presente la prevalencia de cada variable
y la exactitud trivial asociada, recogidas en la
Tabla~\ref{tab:iris-prevalencia-imp} del Capítulo~\ref{sec:methodology}.

\subsection{Resultados del núcleo: \textit{baseline} frente a Agent Skills}
\label{subsec:iris_core}

La Tabla~\ref{tab:iris-b0b1} compara ambos niveles en los cuatro modelos
evaluados: los tres accesibles por API de pago y el modelo local
\texttt{gemma4:e4b}.

\begin{table}[H]
    \centering
    \begingroup
    \scriptsize
    \setlength{\tabcolsep}{4pt}
    \newcommand{\dr}[1]{\textcolor{gray!75}{$#1$}} % fila delta: en gris
    \renewcommand{\texttt}[1]{{\ttfamily #1}}
    \ttabbox[\FBwidth]{
        \caption{Comparación B0 (\textit{baseline}) frente a B1 (Agent Skills). Medias sobre las cinco variables, $N=1.313$. Las filas $\Delta$ (en gris) recogen la variación B1$-$B0.}
        \label{tab:iris-b0b1}
    }{
        \begin{tabularx}{\textwidth}{
            >{\hsize=1.7\hsize\raggedright\arraybackslash}X
            >{\hsize=0.6\hsize\centering\arraybackslash}X
            >{\hsize=0.925\hsize\centering\arraybackslash}X
            >{\hsize=0.925\hsize\columncolor{gray!12}\centering\arraybackslash}X
            >{\hsize=0.925\hsize\centering\arraybackslash}X
            >{\hsize=0.925\hsize\centering\arraybackslash}X
        }
            \toprule
            \textbf{Modelo} & \textbf{Nivel} & \textbf{Exactitud} & \boldmath$\kappa$ & \textbf{F1 macro} & \textbf{Coste (USD)} \\
            \midrule
            \multirow{3}{*}{\texttt{gemini-3.1-flash-lite}}
              & B0 & 0,600 & 0,026 & 0,388 & 5,64 \\
              & B1 & 0,661 & \textbf{0,066} & 0,443 & 13,16 \\
              & \dr{\Delta} & \dr{+0,061} & \dr{+0,040} & \dr{+0,055} & \dr{2{,}3\times} \\
            \midrule
            \multirow{3}{*}{\texttt{gpt-4o-mini}}
              & B0 & 0,564 & 0,033 & 0,410 & 2,52 \\
              & B1 & 0,612 & 0,014 & 0,438 & 6,89 \\
              & \dr{\Delta} & \dr{+0,048} & \dr{-0,019} & \dr{+0,029} & \dr{2{,}7\times} \\
            \midrule
            \multirow{3}{*}{\texttt{gpt-5.4-nano}}
              & B0 & 0,571 & 0,052 & 0,430 & 5,67 \\
              & B1 & 0,569 & 0,004 & 0,364 & 10,86 \\
              & \dr{\Delta} & \dr{-0,003} & \dr{-0,048} & \dr{-0,066} & \dr{1{,}9\times} \\
            \midrule
            \multirow{2}{*}{\textbf{Media (API)}}
              & \textbf{B0} & 0,578 & \textbf{0,037} & 0,409 & \textemdash \\
              & \textbf{B1} & \textbf{0,614} & 0,028 & 0,415 & \textemdash \\
            \midrule
            \multirow{3}{*}{\texttt{gemma4:e4b}\,\textnormal{\footnotesize(local)}}
              & B0 & 0,597 & \textbf{0,061} & 0,421 & \textemdash \\
              & B1 & 0,605 & 0,043 & 0,412 & \textemdash \\
              & \dr{\Delta} & \dr{+0,008} & \dr{-0,018} & \dr{-0,009} & \textemdash \\
            \bottomrule
        \end{tabularx}
    }
    \endgroup
\end{table}

A primera vista, la arquitectura parece funcionar: la exactitud sube con B1 en la
mayoría de los modelos (Figura~\ref{fig:iris-b0b1-metricas}), y la exactitud media de
las API pasa de 0,578 a 0,614. Pero exactitud y $\kappa$ \emph{no cuentan la
misma historia}: el $\kappa$ medio, lejos de acompañar, baja de 0,037 a 0,028, y todo
ello a un coste de inferencia 2,2 veces mayor. Esta discordancia obliga a precisar
qué métrica resulta pertinente antes de interpretar el efecto de las \textit{skills}:
en una tarea con clases fuertemente desbalanceadas, la exactitud es un indicador
engañoso y el análisis debe apoyarse en el coeficiente $\kappa$, según se argumenta
en la Sección~\ref{subsec:iris_divergencia}.

\begin{figure}[H]
    \centering
    \includegraphics[width=\textwidth]{fig_b0b1_barras_blanco.pdf}
    \caption{Niveles de exactitud, F1 macro y $\kappa$ por modelo, en B0 y B1
    ($N=1.313$). Cada métrica se representa en su propia escala: el $\kappa$ vive
    en un orden de magnitud mucho menor que la exactitud.}
    \label{fig:iris-b0b1-metricas}
\end{figure}


\subsection{Divergencia entre exactitud y \texorpdfstring{$\kappa$}{kappa}}
\label{subsec:iris_divergencia}

Conviene detenerse en esa discordancia, porque condiciona cómo debe leerse toda la
tabla: \textbf{la mejora en exactitud no se traduce en más acuerdo}. Considerando
las medias, el nivel B1 sube la exactitud en 0,036 pero el $\kappa$ baja en $-0,009$.
Exactitud y $\kappa$ pueden, por tanto, moverse de forma independiente e incluso en
sentidos opuestos. Esto es un efecto conocido de $\kappa$ en
presencia de clases desbalanceadas~\cite{feinstein1990}, pero aquí resulta determinante: pasarlo por alto
invertiría la conclusión del capítulo. La Figura~\ref{fig:iris-divergencia} lo
visualiza como trayectorias B0$\to$B1: en el panel de exactitud casi todas las líneas
ascienden, mientras que en el de $\kappa$ se mantienen o descienden, salvo en
\texttt{gemini}, cuya línea sube en ambas.

\begin{figure}[H]
    \centering
    \includegraphics[width=\textwidth]{fig_b0b1_slope_blanco.pdf}
    \caption{Evolución de cada modelo de B0 a B1 en las tres métricas ($N=1.313$).
    La exactitud sube en tres de los cuatro modelos, pero el $\kappa$ solo mejora en
    \texttt{gemini-3.1-flash-lite}. En el resto se mantiene o desciende. La mejora en
    exactitud sin mejora en $\kappa$ es el desacoplamiento entre exactitud y acuerdo.}
    \label{fig:iris-divergencia}
\end{figure}

El motivo de esa divergencia es conocido. Al volverse menos conservador, el modelo
responde \textit{Sí} con más frecuencia, y cuando la clase positiva es la mayoritaria,
contestar afirmativamente más a menudo eleva la exactitud por simple coincidencia
con esa mayoría, no por discriminar mejor. El coeficiente $\kappa$, que descuenta el
acuerdo esperable por azar, no premia ese desplazamiento. La consecuencia práctica
es directa: \emph{un estudio que reportara solo exactitud concluiría que las Agent
Skills mejoran el sistema, y se equivocaría}. De ahí la recomendación de reportar
siempre exactitud, $\kappa$ y prevalencia de forma conjunta en tareas de
clasificación desbalanceadas.

Leído correctamente, por $\kappa$, el efecto de las \textit{Agent Skills}
\textbf{depende del modelo}. Solo \texttt{gemini-3.1-flash-lite} mejora con la
arquitectura: su $\kappa$ pasa de 0,026 a 0,066, y lo hace acompañado de subidas
simultáneas en exactitud y F1. En los dos modelos de OpenAI las \textit{skills} no ayudan e incluso
perjudican ($\Delta\kappa = -0,019$ en \texttt{gpt-4o-mini} y $-0,048$ en
\texttt{gpt-5.4-nano}), y el modelo local se comporta igual que ellos ($\kappa$ de
0,061 a 0,043). De los cuatro modelos, solo uno mejora con las \textit{skills}. En
los otros tres el $\kappa$ se mantiene o desciende, y la media no aumenta. Que el
efecto dependa del modelo y no eleve el acuerdo en agregado admite varias lecturas, desde el techo de la propia tarea hasta la heterogeneidad de la anotación, que se
retoman en la discusión (Sección~\ref{sec:discussion}) una vez presentada toda la
evidencia.

\subsection{Descomposición de la arquitectura por componentes}
\label{subsec:iris_descomposicion}

El nivel B1 combina tres ingredientes, descritos en el Capítulo~\ref{implementation}:
la carga de \textit{skills} bajo demanda con verificación de evidencias
(Sección~\ref{subsec:diseno_agentico}), un catálogo de resúmenes de las guías
expertas descubrible por el agente (Sección~\ref{prompts}), y una herramienta de
recuperación en vivo sobre esas guías (RAG, implementada como \texttt{CONSULTAR\_GUIA},
Sección~\ref{rag}). Para averiguar cuál de ellos explica el efecto
observado se ejecutaron cinco variantes por modelo, activando y desactivando cada
componente de forma independiente, sobre los cuatro modelos y las 1.313 noticias.
Las Tablas~\ref{tab:iris-descomposicion}, \ref{tab:iris-descomposicion-acc}
y~\ref{tab:iris-descomposicion-f1} recogen, respectivamente, el $\kappa$, la
exactitud y el F1 macro medios de cada combinación.

\begin{table}[H]
    \centering
    \begingroup
    \scriptsize
    \ttabbox[\FBwidth]{
        \caption{Descomposición de B1 en componentes. $\kappa$ medio sobre las cinco variables, por modelo ($N=1.313$)}
        \label{tab:iris-descomposicion}
    }{
        \begin{tabularx}{\textwidth}{
            X
            >{\centering\arraybackslash}X
            >{\centering\arraybackslash}X
            >{\centering\arraybackslash}X
            >{\centering\arraybackslash}X
        }
            \toprule
            \textbf{Configuración} & \textbf{Gemini} & \textbf{GPT-4o-mini} & \textbf{GPT-5.4-nano} & \textbf{Gemma} \\
            \midrule
            B0 & 0,026 & 0,033 & 0,052 & 0,061 \\
            Solo \textit{skills} & 0,061 & 0,025 & 0,000 & 0,065 \\
            \textit{Skills} + resúmenes & \textbf{0,075} & 0,024 & 0,015 & 0,061 \\
            \textit{Skills} + RAG & 0,063 & 0,028 & 0,012 & 0,038 \\
            Completo & 0,066 & 0,014 & 0,004 & 0,043 \\
            \bottomrule
        \end{tabularx}
    }
    \endgroup
\end{table}

\begin{table}[H]
    \centering
    \begingroup
    \scriptsize
    \ttabbox[\FBwidth]{
        \caption{Descomposición de B1 en componentes. \textbf{Exactitud} media sobre las cinco variables, por modelo ($N=1.313$)}
        \label{tab:iris-descomposicion-acc}
    }{
        \begin{tabularx}{\textwidth}{
            X
            >{\centering\arraybackslash}X
            >{\centering\arraybackslash}X
            >{\centering\arraybackslash}X
            >{\centering\arraybackslash}X
        }
            \toprule
            \textbf{Configuración} & \textbf{Gemini} & \textbf{GPT-4o-mini} & \textbf{GPT-5.4-nano} & \textbf{Gemma} \\
            \midrule
            B0 & 0,600 & 0,564 & 0,571 & 0,597 \\
            Solo \textit{skills} & 0,645 & \textbf{0,677} & 0,567 & 0,608 \\
            \textit{Skills} + resúmenes & 0,656 & 0,609 & 0,569 & 0,610 \\
            \textit{Skills} + RAG & 0,651 & 0,631 & 0,571 & 0,602 \\
            Completo & 0,661 & 0,612 & 0,569 & 0,605 \\
            \bottomrule
        \end{tabularx}
    }
    \endgroup
\end{table}

\begin{table}[H]
    \centering
    \begingroup
    \scriptsize
    \ttabbox[\FBwidth]{
        \caption{Descomposición de B1 en componentes. \textbf{F1 macro} media sobre las cinco variables, por modelo ($N=1.313$)}
        \label{tab:iris-descomposicion-f1}
    }{
        \begin{tabularx}{\textwidth}{
            X
            >{\centering\arraybackslash}X
            >{\centering\arraybackslash}X
            >{\centering\arraybackslash}X
            >{\centering\arraybackslash}X
        }
            \toprule
            \textbf{Configuración} & \textbf{Gemini} & \textbf{GPT-4o-mini} & \textbf{GPT-5.4-nano} & \textbf{Gemma} \\
            \midrule
            B0 & 0,388 & 0,410 & 0,430 & 0,421 \\
            Solo \textit{skills} & 0,430 & \textbf{0,484} & 0,364 & 0,424 \\
            \textit{Skills} + resúmenes & 0,444 & 0,451 & 0,365 & 0,423 \\
            \textit{Skills} + RAG & 0,439 & 0,473 & 0,374 & 0,409 \\
            Completo & 0,443 & 0,438 & 0,364 & 0,412 \\
            \bottomrule
        \end{tabularx}
    }
    \endgroup
\end{table}

\begin{figure}[H]
    \centering
    \includegraphics[width=\textwidth]{fig_kappa_componentes_blanco.pdf}
    \caption{Coeficiente $\kappa$ medio al activar cada componente de B1, por modelo
    ($N=1.313$). La mejor configuración de cada modelo prescinde de la recuperación
    en vivo (RAG, franja sombreada). El mecanismo de \textit{skills} solo eleva el
    acuerdo en \texttt{gemini-3.1-flash-lite}.}
    \label{fig:iris-descomposicion-fig}
\end{figure}

Dos conclusiones se sostienen en los cuatro modelos, y la
Figura~\ref{fig:iris-descomposicion-fig} las resume. La primera afecta a la
herramienta RAG de recuperación en vivo: no aporta, y con frecuencia perjudica. En
todos los modelos la mejor variante prescinde de ella. En \texttt{gemma4:e4b} el
contraste es nítido: sin la RAG el $\kappa$ se sitúa entre 0,061 y
0,065, y al activarla cae a valores entre 0,038 y 0,043. La recuperación
por TF-IDF sobre textos periodísticos cortos en este caso introduce más ruido que señal.

La segunda conclusión concierne al mecanismo de \textit{skills}: solo mejora el
acuerdo en \texttt{gemini}. En ese modelo, la variante que se limita a cargar la
metodología como \textit{skill} verificable, sin guías de ningún tipo, ya recupera
la mayor parte de la mejora sobre el control, de 0,026 a 0,061. En los otros
tres modelos esa misma variante iguala el control (\texttt{gemma}) o lo empeora
(los dos de OpenAI). Cuando hay ganancia, la configuración con mejor acuerdo es
\textit{skills} más resúmenes sin RAG: más simple y más barata que el
sistema completo, y en ningún modelo peor.

En conjunto, la descomposición confirma la lectura del núcleo: la arquitectura no es
útil ni inútil en abstracto, sino que \textbf{amplifica lo que cada modelo ya trae}.
Aporta cuando el modelo aprovecha de verdad la carga bajo demanda, y estorba cuando la
infrautiliza o se satura con el catálogo. Su valor está condicionado a la capacidad
del modelo para explotarla.

El desglose en exactitud y F1 macro matiza el cuadro y, de paso, ilustra la
divergencia de la Sección~\ref{subsec:iris_divergencia}: la exactitud puede subir con
las \textit{skills} (\texttt{gpt-4o-mini} con solo \textit{skills} pasa de 0,564 a
0,677) sin que el acuerdo la acompañe (su $\kappa$ se queda en 0,025).


\subsection{Diagnóstico por variable}
\label{subsec:iris_diagnostico}

Hasta ahora el análisis se ha apoyado en métricas globales, promediadas sobre las
cinco variables. Ese nivel agregado responde a la pregunta central del capítulo,
pero esconde comportamientos muy distintos de una variable a otra. Para entender qué
ocurre por debajo del promedio hay que desglosar los resultados variable a variable.
Para cada una se reportan, sobre la clase positiva (\textit{Sí}), la \textbf{precisión} (de
las piezas que el modelo marca como \textit{Sí}, cuántas lo son de verdad, que no debe
confundirse con la exactitud, la cual mide el acierto global sobre todas las piezas)
y el \textbf{recall} o sensibilidad (del sexismo que marcan las expertas, cuánto
detecta el modelo), junto al coeficiente $\kappa$. El recall es aquí el más revelador: un recall bajo significa
que el modelo deja sin marcar buena parte de lo que las anotadoras sí marcan. La
Tabla~\ref{tab:iris-diagnostico} recoge estos valores para el nivel B1, combinándolos con la
revisión manual de casos se caracteriza cada variable.

\begin{table}[H]
    \centering
    \begingroup
    \scriptsize
    \setlength{\tabcolsep}{3pt}
    \newcommand{\tri}[3]{\begin{tabular}[t]{@{}c@{}}#1\\#2\\#3\end{tabular}}
    \ttabbox[\FBwidth]{
        \caption{Precisión, recall (sensibilidad) sobre la clase \textit{Sí} y $\kappa$, por variable en el nivel B1 ($N=1.313$). En cada celda se apilan, de arriba abajo, precisión, recall y $\kappa$. ``Prev.'' es la prevalencia de la clase \textit{Sí}.}
        \label{tab:iris-diagnostico}
    }{
        \begin{tabularx}{\textwidth}{
            >{\hsize=1.35\hsize\raggedright\arraybackslash\ttfamily}X
            >{\hsize=0.55\hsize\centering\arraybackslash}X
            >{\hsize=0.775\hsize\centering\arraybackslash}X
            >{\hsize=0.775\hsize\centering\arraybackslash}X
            >{\hsize=0.775\hsize\centering\arraybackslash}X
            >{\hsize=0.775\hsize\centering\arraybackslash}X
        }
            \toprule
            \normalfont\textbf{Variable} & \textbf{Prev.} & \textbf{Gemini} & \textbf{GPT-4o-mini} & \textbf{GPT-5.4-nano} & \textbf{Gemma} \\
            \midrule
            lenguaje\_sexista & 0,810 & \tri{0,961}{0,046}{0,015} & \tri{0,789}{0,450}{$-$0,037} & \tri{1,000}{0,001}{0,000} & \tri{0,929}{0,012}{0,003} \\
            \midrule
            masc\_generico & 0,810 & \textbf{\tri{0,887}{0,733}{0,261}} & \tri{0,873}{0,206}{0,034} & \tri{0,874}{0,104}{0,016} & \tri{0,897}{0,408}{0,109} \\
            \midrule
            sexismo\_discurso & 0,428 & \tri{0,423}{0,020}{0,000} & \tri{0,450}{0,112}{0,011} & \tri{0,435}{0,114}{0,004} & \tri{0,378}{0,025}{$-$0,006} \\
            \midrule
            asimetria\_mujer\_hombre & 0,062 & \tri{0,118}{0,074}{0,045} & \tri{0,063}{0,148}{0,003} & \tri{0,000}{0,000}{$-$0,012} & \tri{0,143}{0,160}{0,092} \\
            \midrule
            denominacion\_sexualizada & 0,100 & \tri{0,167}{0,015}{0,011} & \tri{0,375}{0,046}{0,061} & \tri{1,000}{0,008}{0,014} & \tri{0,158}{0,023}{0,015} \\
            \bottomrule
        \end{tabularx}
    }
    \endgroup
\end{table}

\begin{itemize}
    \item \textbf{\texttt{lenguaje\_sexista} (V25): un criterio en disputa.} La
    precisión es alta cuando el modelo marca \textit{Sí} (entre 0,79 y 1,00), pero el
    $\kappa$ ronda el azar en los cuatro modelos (entre $-0,037$ y 0,015). El problema
    no es que el modelo se equivoque al marcar, sino que sistema y anotadoras no
    coinciden en \emph{cuándo} hay que marcar. La causa está en el propio criterio de
    V25, que es intrínsecamente amplio y admite dos lecturas. Al ser la variable más
    general del bloque (el uso discriminatorio del lenguaje por razón de sexo, sin
    acotar a un mecanismo concreto), las anotadoras se inclinan por una lectura
    inclusiva y la marcan en casi todas las piezas, de ahí su prevalencia del
    81\,\%, mientras que el modelo se atiene a una lectura más literal y la marca
    mucho menos, de ahí su recall bajísimo. Ninguna de las dos lecturas es
    incorrecta, porque el propio criterio admite ambas.

    \item \textbf{\texttt{masc\_generico} (V26): el único acuerdo genuino, y su
    umbral.} Es el único caso con acuerdo apreciable: \texttt{gemini} alcanza
    $\kappa = 0,261$ y un recall de 0,733, es decir, detecta la mayoría del masculino
    genérico que marcan las expertas, con exactitud y $\kappa$ que mejoran a la vez
    sobre el \textit{baseline} (lo que descarta el artefacto de marginales y confirma
    una ganancia real). El local \texttt{gemma4:e4b} queda en segundo lugar, con un
    recall de 0,408 y un $\kappa$ de 0,109, la única variable en la que también supera
    con claridad el azar. Los dos modelos de OpenAI, en cambio, tienen un recall mucho
    menor (0,206 y 0,104): reflejan un desajuste de umbral, porque el español emplea el
    masculino genérico de forma pervasiva (\textit{los trabajadores}, \textit{los
    ciudadanos}) y la anotación tiende a marcarlo de forma sistemática, mientras que
    los modelos solo lo señalan cuando resulta evitable e invisibilizador.

    \item \textbf{\texttt{sexismo\_discurso} (V30): anotación sobre-inclusiva.} El
    acuerdo ronda el azar en los cuatro modelos ($\kappa$ entre $-0,006$ y 0,011) y el
    recall es muy bajo (entre 0,020 y 0,114). En una fracción de los casos positivos
    no se aprecia, en una lectura atenta del texto, el fenómeno codificado. Estos
    casos se concentran en el anotador de menor acuerdo global (véase la
    Sección~\ref{subsec:iris_errores}).

    \item \textbf{\texttt{asimetria\_mujer\_hombre} (V33) y
    \texttt{denominacion\_sexualizada} (V35): desbalance extremo.} Con prevalencias
    del 6,2\,\% y el 10,0\,\%, el modelo casi no marca positivos (recall entre 0,000 y
    0,160) y la precisión, cuando no es nula, descansa sobre un puñado de
    predicciones: \texttt{gpt-5.4-nano} marca \textit{Sí} en \texttt{denominacion\_sexualizada}
    tan pocas veces que su precisión de 1,000 carece de valor. En este régimen
    $\kappa$ penaliza severamente incluso a clasificadores razonables, hasta el punto
    de que una exactitud de 0,90 puede coexistir con $\kappa = 0,000$.
\end{itemize}

El desglose por variable revela además dónde actúan las \textit{skills}. La
Tabla~\ref{tab:iris-porvariable-b0b1} compara el $\kappa$ de B0 y B1 en cada
variable. En \texttt{gemini}, la mejora del núcleo se concentra por completo en
\texttt{masc\_generico} (de 0,073 a 0,261), mientras el resto de variables apenas se
mueve. En los otros tres modelos ocurre lo contrario: las \textit{skills} degradan
variables donde el control ya tenía algo de señal, sobre todo
\texttt{denominacion\_sexualizada} (de 0,123 a 0,014 en \texttt{gpt-5.4-nano}, de
0,136 a 0,061 en \texttt{gpt-4o-mini} y de 0,154 a 0,015 en \texttt{gemma4:e4b}). El
efecto de la arquitectura no solo depende del modelo, sino que dentro de cada modelo
se localiza en variables concretas.

\begin{table}[H]
    \centering
    \begingroup
    \scriptsize
    \ttabbox[\FBwidth]{
        \caption{Coeficiente $\kappa$ por variable en B0 (\textit{baseline}) y B1 (Agent Skills), para los cuatro modelos ($N=1.313$)}
        \label{tab:iris-porvariable-b0b1}
    }{
        \begin{tabularx}{\textwidth}{
            >{\hsize=1.4\hsize\raggedright\arraybackslash\ttfamily}X
            >{\hsize=0.9\hsize\centering\arraybackslash}X
            >{\hsize=0.9\hsize\centering\arraybackslash}X
            >{\hsize=0.9\hsize\centering\arraybackslash}X
            >{\hsize=0.9\hsize\centering\arraybackslash}X
        }
            \toprule
            \normalfont\textbf{Variable} & \textbf{Gemini} & \textbf{GPT-4o-mini} & \textbf{GPT-5.4-nano} & \textbf{Gemma} \\
             & \multicolumn{4}{c}{\normalfont\scriptsize $\kappa$: B0 / B1} \\
            \midrule
            lenguaje\_sexista & 0,000 / 0,015 & $-0,013$ / $-0,037$ & 0,001 / 0,000 & $-0,004$ / 0,003 \\
            masc\_generico & \textbf{0,073 / 0,261} & 0,023 / 0,034 & 0,103 / 0,016 & 0,102 / 0,109 \\
            sexismo\_discurso & 0,004 / 0,000 & 0,002 / 0,011 & 0,007 / 0,004 & 0,002 / $-0,006$ \\
            asimetria\_mujer\_hombre & 0,036 / 0,045 & 0,019 / 0,003 & 0,026 / $-0,012$ & 0,054 / 0,092 \\
            denominacion\_sexualizada & 0,019 / 0,011 & 0,136 / 0,061 & 0,123 / 0,014 & 0,154 / 0,015 \\
            \bottomrule
        \end{tabularx}
    }
    \endgroup
\end{table}

Conviene mirar también cuánto marca cada modelo en términos absolutos, que es el
reverso del recall: marcar poco implica dejar sin detectar buena parte de lo que las
expertas sí señalan. La Tabla~\ref{tab:iris-prev-modelo-var} recoge, por variable, la
prevalencia de \textit{Sí} de cada modelo. La única variable en la que algún modelo
se acerca al terreno de las expertas es \texttt{masc\_generico} (\texttt{gemini},
66,9\,\%), mientras que en el resto todos marcan \textit{Sí} en una fracción muy
pequeña de las piezas. El patrón, además, no es común: \texttt{gpt-4o-mini} dispara
en \texttt{lenguaje\_sexista} (44,6\,\%) donde \texttt{gemini} apenas marca (3,8\,\%),
mientras que \texttt{gemini} domina en \texttt{masc\_generico}. Cada modelo tiene su
propio umbral por variable.

\begin{table}[H]
    \centering
    \begingroup
    \scriptsize
    \renewcommand{\texttt}[1]{{\ttfamily #1}}
    \ttabbox[\FBwidth]{
        \caption{Prevalencia de \textit{Sí} (\%) de cada modelo por variable, en el nivel B1 ($N=1.313$)}
        \label{tab:iris-prev-modelo-var}
    }{
        \begin{tabularx}{\textwidth}{
            >{\hsize=1.6\hsize\raggedright\arraybackslash\ttfamily}X
            >{\hsize=0.85\hsize\centering\arraybackslash}X
            >{\hsize=0.85\hsize\centering\arraybackslash}X
            >{\hsize=0.85\hsize\centering\arraybackslash}X
            >{\hsize=0.85\hsize\centering\arraybackslash}X
        }
            \toprule
            \normalfont\textbf{Variable} & \textbf{Gemini} & \textbf{GPT-4o-mini} & \textbf{GPT-5.4-nano} & \textbf{Gemma} \\
            \midrule
            lenguaje\_sexista & 3,8 & 44,6 & 0,1 & 1,1 \\
            masc\_generico & \textbf{66,9} & 19,1 & 9,7 & 36,9 \\
            sexismo\_discurso & 2,0 & 10,7 & 11,2 & 2,8 \\
            asimetria\_mujer\_hombre & 3,9 & 14,4 & 0,7 & 6,9 \\
            denominacion\_sexualizada & 0,9 & 1,2 & 0,1 & 1,4 \\
            \bottomrule
        \end{tabularx}
    }
    \endgroup
\end{table}

Un rasgo atraviesa a las cinco variables: el recall es bajo casi en todas, es decir,
los modelos detectan solo una fracción del sexismo que marcan las expertas. Esta
\textbf{infradetección generalizada} es el patrón dominante del capítulo, y se
examina de frente en la Sección~\ref{subsec:iris_errores} para dirimir si es una
limitación de los modelos o de la propia tarea.

\subsection{El modelo local: rendimiento y coste}
\label{subsec:iris_local}\label{subsec:iris_coste}

Entre los cuatro modelos, el local \texttt{gemma4:e4b} merece una lectura aparte. En
el nivel de control (B0) obtiene el \textbf{mayor $\kappa$ de los cuatro} (0,061), por
encima de las tres API de pago (0,026, 0,033 y 0,052), y ante las \textit{skills} se
comporta como los modelos de OpenAI: su acuerdo baja de 0,061 a 0,043
(Tabla~\ref{tab:iris-descomposicion}). Que un modelo abierto, gratuito y desplegable
en infraestructura propia iguale, y en $\kappa$ supere, a los servicios
comerciales tiene consecuencias prácticas para un proyecto que maneja material
sujeto a privacidad. Ahora bien, esa ventaja en acuerdo y en ausencia de tarifas
convive con un coste de otro tipo: el del cómputo.

La Tabla~\ref{tab:iris-b0b1} recoge el coste real de inferencia. La arquitectura de
agentes multiplica por 2,2 el coste en API respecto al \textit{baseline} (30,91 USD
frente a 13,83 USD para los tres modelos de pago sobre 1.313 noticias): el agente
encadena en promedio entre dos y tres llamadas por variable (una por cada acción del
bucle más el veredicto final) frente a la única llamada del control. La diferencia no
es mayor porque las llamadas de B0 son más largas, al inyectar la metodología
completa en el \textit{prompt}.

En el modelo local no hay tarifas, pero el coste reaparece como tiempo de cómputo. El
nivel B1 de \texttt{gemma4:e4b} se ejecutó sobre la granja de GPU (NVIDIA RTX 3090 y
4090) del Departamento de Teoría de la Señal y Comunicaciones (TSC) de la Universidad
Carlos III de Madrid, con un rendimiento medido de unos 44 segundos por variable,
esto es, en torno a 220 segundos por artículo completo. Repartiendo la carga en
cuatro GPU, cada corrida sobre el corpus de evaluación tardó alrededor de veinte
horas. En una sola GPU, el mismo cómputo superaría las ochenta horas. Así, el modelo
local es competitivo en acuerdo y prescinde de tarifas, pero su despliegue agéntico
exige infraestructura dedicada: el coste no desaparece, cambia de dinero a cómputo.
La Figura~\ref{fig:iris-acuerdo-dolar} resume ese compromiso: situado el acuerdo del
sistema desplegado (B1) frente a su coste monetario, el modelo local ocupa la esquina
más favorable (máximo acuerdo por dólar), con la salvedad de que su coste se traslada
al cómputo. Las implicaciones prácticas para IRIS se retoman en la discusión
(Sección~\ref{sec:discussion}).

\begin{figure}[H]
    \centering
    \includegraphics[width=\textwidth]{fig_acuerdo_por_dolar_blanco.pdf}
    \caption{Acuerdo (coeficiente $\kappa$ en B1) frente al coste de inferencia sobre
    el corpus completo, por modelo ($N=1.313$). El modelo local ofrece el mejor
    acuerdo por dólar (coste monetario nulo). Su coste se traslada al tiempo de
    cómputo en infraestructura propia.}
    \label{fig:iris-acuerdo-dolar}
\end{figure}

\subsection{¿Un límite del modelo o de la tarea?}
\label{subsec:iris_errores}

El bajo acuerdo absoluto plantea la pregunta que condiciona la interpretación de
todo el capítulo: ¿refleja una limitación de los modelos, o un límite intrínseco de
la propia tarea? Tres análisis complementarios apuntan en la misma dirección.

\subsubsection{Heterogeneidad entre equipos de anotación}

El corpus registra, para cada pieza, la persona que la codificó. Las 1.313
noticias anotadas se reparten entre dos equipos: Indexa (548 piezas) y UCM3
(765). Aunque la ausencia de doble codificación
(Sección~\ref{subsec:limitacion_fiabilidad}) impide calcular la fiabilidad
intercodificadora, sí es posible medir el acuerdo del sistema con cada equipo
por separado, a modo de análisis de sensibilidad. La
Tabla~\ref{tab:iris-sensibilidad} recoge el resultado para el nivel B1.

\begin{table}[H]
    \centering
    \begingroup
    \scriptsize
    \renewcommand{\texttt}[1]{{\ttfamily #1}}
    \ttabbox[\FBwidth]{
        \caption{Análisis de sensibilidad: acuerdo (B1) según el equipo de anotación}
        \label{tab:iris-sensibilidad}
    }{
        \begin{tabularx}{\textwidth}{
            >{\hsize=1.5\hsize\raggedright\arraybackslash}X
            >{\hsize=1.0\hsize\raggedright\arraybackslash}X
            >{\hsize=0.85\hsize\centering\arraybackslash}X
            >{\hsize=0.85\hsize\centering\arraybackslash}X
            >{\hsize=0.8\hsize\columncolor{gray!15}\centering\arraybackslash}X
        }
            \toprule
            \textbf{Modelo} & \textbf{Subconjunto} & \textbf{Exactitud} & \textbf{F1 macro} & \boldmath$\kappa$ \\
            \midrule
            \multirow{3}{*}{\texttt{gemini-3.1-flash-lite}}
              & Todas (1.313)   & 0,661 & 0,443 & 0,066 \\
              & Solo Indexa (548) & \textbf{0,741} & \textbf{0,503} & \textbf{0,118} \\
              & Solo UCM3 (765)  & 0,604 & 0,398 & 0,043 \\
            \midrule
            \multirow{3}{*}{\texttt{gpt-4o-mini}}
              & Todas           & 0,612 & 0,438 & 0,014 \\
              & Solo Indexa     & 0,693 & 0,482 & 0,042 \\
              & Solo UCM3       & 0,554 & 0,398 & 0,001 \\
            \midrule
            \multirow{3}{*}{\texttt{gpt-5.4-nano}}
              & Todas           & 0,569 & 0,364 & 0,004 \\
              & Solo Indexa     & 0,661 & 0,406 & 0,014 \\
              & Solo UCM3       & 0,502 & 0,323 & $-0,001$ \\
            \midrule
            \multirow{3}{*}{\texttt{gemma4:e4b}}
              & Todas           & 0,605 & 0,411 & 0,043 \\
              & Solo Indexa     & 0,695 & 0,461 & 0,068 \\
              & Solo UCM3       & 0,541 & 0,369 & 0,032 \\
            \bottomrule
        \end{tabularx}
    }
    \endgroup
\end{table}

El acuerdo con Indexa es sistemáticamente muy superior al obtenido con UCM3, en
los cuatro modelos y en las tres métricas de forma coherente. Para
\texttt{gemini-3.1-flash-lite}, restringir la evaluación a Indexa prácticamente
duplica el coeficiente $\kappa$ (de 0,066 a 0,118) y eleva la exactitud
ocho puntos. Contra UCM3 en solitario, $\kappa$ desciende hacia el azar (hasta
valores negativos en \texttt{gpt-5.4-nano}). \textbf{La magnitud de esta
variabilidad entre equipos es comparable a la diferencia entre modelos}, lo que
indica que el origen de la anotación pesa en el acuerdo tanto como la elección
del modelo.

El origen de esa diferencia está en la anotación, no en el modelo. La
Tabla~\ref{tab:iris-equipos-var} muestra, por variable, la proporción de \textit{Sí}
que marca cada equipo, y deja ver que la misma variable se codifica de forma muy
distinta. El caso extremo es \texttt{sexismo\_discurso}: Indexa marca el fenómeno en
el 11,5\,\% de las piezas y UCM3 en el 65,2\,\%, casi seis veces más. No es ruido del
modelo, son dos criterios humanos difícilmente conciliables sobre una misma variable.
Como esa variabilidad entre equipos iguala o supera la que separa a los modelos,
constituye evidencia directa del techo de la tarea, sin necesitar doble codificación.

\begin{table}[H]
    \centering
    \begingroup
    \scriptsize
    \renewcommand{\texttt}[1]{{\ttfamily #1}}
    \ttabbox[\FBwidth]{
        \caption{Prevalencia de \textit{Sí} y acuerdo (\texttt{gemini} B1) por equipo y variable. La misma variable se codifica de forma muy distinta según el equipo.}
        \label{tab:iris-equipos-var}
    }{
        \begin{tabularx}{\textwidth}{
            >{\hsize=1.7\hsize\raggedright\arraybackslash\ttfamily}X
            >{\hsize=0.9\hsize\centering\arraybackslash}X
            >{\hsize=0.9\hsize\centering\arraybackslash}X
            >{\hsize=0.9\hsize\centering\arraybackslash}X
            >{\hsize=0.8\hsize\centering\arraybackslash}X
            >{\hsize=0.8\hsize\centering\arraybackslash}X
        }
            \toprule
             & \multicolumn{2}{c}{\normalfont\% \textit{Sí}} & & \multicolumn{2}{c}{\normalfont$\kappa$} \\
            \cmidrule(lr){2-3}\cmidrule(lr){5-6}
            \normalfont\textbf{Variable} & \textbf{Indexa} & \textbf{UCM3} & \textbf{Ratio} & \textbf{Indexa} & \textbf{UCM3} \\
            \midrule
            lenguaje\_sexista & 74,6 & 85,5 & 1,1$\times$ & 0,026 & 0,009 \\
            masc\_generico & 71,4 & 88,0 & 1,2$\times$ & 0,277 & 0,223 \\
            sexismo\_discurso & \textbf{11,5} & \textbf{65,2} & \textbf{5,7$\times$} & 0,065 & $-0,003$ \\
            asimetria\_mujer\_hombre & 3,3 & 8,2 & 2,5$\times$ & 0,202 & $-0,027$ \\
            denominacion\_sexualizada & 6,6 & 12,4 & 1,9$\times$ & 0,019 & 0,013 \\
            \bottomrule
        \end{tabularx}
    }
    \endgroup
\end{table}

Debe subrayarse que este resultado \emph{no} autoriza a reportar únicamente el
subconjunto más favorable: al no existir doble codificación, no puede
determinarse qué equipo refleja mejor el constructo, y la diferencia está además
confundida con el hecho de que cada equipo codificó piezas distintas. El
análisis acota el techo alcanzable, pero no lo atribuye a un error de anotación.

\subsubsection{Consenso de los modelos frente al \textit{ground truth}}

Un segundo análisis examina las celdas en las que \textbf{los cuatro modelos
coinciden entre sí} pero difieren de la anotación experta: son las más
diagnósticas, pues descartan un error achacable a un solo modelo. De
las 6.565 celdas evaluadas (1.313 piezas $\times$ 5 variables), 1.364 presentan
consenso de los cuatro modelos en contra del \textit{ground truth}. Su dirección
es prácticamente unánime: en \textbf{1.363 casos} los cuatro modelos responden \textit{No}
donde la anotación dice \textit{Sí}, y solo en 1 caso ocurre lo contrario. La
infradetección observada por variable es, por tanto, un \textbf{sesgo sistemático},
no aleatorio: los modelos detectan de forma consistente menos sexismo que las
anotadoras.


\subsubsection{El sistema como una anotadora más}
\label{subsec:iris_equipo_ia}

Los análisis anteriores comparan cada modelo contra la anotación experta tomada
como referencia. Un enfoque complementario invierte la pregunta: en vez de medir
cuánto se equivoca el modelo respecto a un patrón fijo, lo sitúa como \textbf{una
anotadora más} dentro del espacio de criterios de las personas que codificaron el
corpus. La cuestión deja de ser ``¿acierta el modelo?'' y pasa a ser ``¿dónde cae su
criterio respecto al de las expertas?''. El sistema se trata así como un tercer
equipo de anotación, junto a Indexa y UCM3.

Como cada noticia la codificó una sola persona, no hay solape humano-humano. Lo que
sí compara a todos sin solape es la \textbf{prevalencia de \textit{Sí}}: la proporción de
piezas en que cada codificador, humano o modelo, marca la variable. La
Tabla~\ref{tab:iris-criterios} ordena por esa proporción, promediada sobre las cinco
variables, a los ocho anotadores humanos y a los cuatro modelos evaluados.

\begin{table}[H]
    \centering
    \begingroup
    \scriptsize
    \renewcommand{\texttt}[1]{{\ttfamily #1}}
    \ttabbox[\FBwidth]{
        \caption{Espacio de criterios: prevalencia media de \textit{Sí} por anotador (humano y modelo)}
        \label{tab:iris-criterios}
    }{
        \begin{tabularx}{\textwidth}{
            >{\hsize=1.5\hsize\raggedright\arraybackslash}X
            >{\hsize=0.8\hsize\centering\arraybackslash}X
            >{\hsize=0.7\hsize\centering\arraybackslash}X
            >{\hsize=1.0\hsize\columncolor{gray!15}\centering\arraybackslash}X
        }
            \toprule
            \textbf{Anotador} & \textbf{Equipo} & \boldmath$n$ & \textbf{Prev. \textit{Sí} media} \\
            \midrule
            UCM3 3 & UCM3 & 363 & 0,639 \\
            UCM3 4 & UCM3 & 166 & 0,465 \\
            Indexa 1 & Indexa & 127 & 0,408 \\
            Indexa 5 & Indexa & 50 & 0,396 \\
            UCM3 5 & UCM3 & 235 & 0,371 \\
            Indexa 2 & Indexa & 128 & 0,325 \\
            Indexa 3 & Indexa & 231 & 0,289 \\
            Indexa 4 & Indexa & 12 & 0,283 \\
            \midrule[\heavyrulewidth]
            \texttt{gpt-4o-mini} & IA & 1.313 & 0,180 \\
            \texttt{gemini-3.1-flash-lite} & IA & 1.313 & 0,155 \\
            \texttt{gemma4:e4b} & IA & 1.313 & 0,099 \\
            \texttt{gpt-5.4-nano} & IA & 1.313 & 0,043 \\
            \bottomrule
        \end{tabularx}
    }
    \endgroup
\end{table}

Los cuatro modelos ocupan el extremo permisivo de la escala, por debajo de la
anotadora humana más contenida. No se sitúan entre Indexa y UCM3: caen fuera del
rango humano completo. Es la infradetección sistemática de la subsección anterior,
leída ahora sobre una escala común a personas y máquinas. Como cada persona
codificó piezas distintas, su prevalencia mezcla dos factores: el criterio propio y
la muestra que le tocó. Para separarlos, se aplicó cada modelo, como instrumento
común, a las mismas piezas que vio cada anotadora. El orden apenas cambia: UCM3-3
sigue siendo, con diferencia, la anotadora más permisiva.

Un segundo análisis mide el acuerdo \emph{entre} modelos, que sí tienen solape
total porque los cuatro codificaron las 1.313 piezas. El resultado es inesperado y
se recoge en la Tabla~\ref{tab:iris-intraia}.

\begin{table}[H]
    \centering
    \begingroup
    \scriptsize
    \renewcommand{\texttt}[1]{{\ttfamily #1}}
    \ttabbox[\FBwidth]{
        \caption{Acuerdo entre modelos ($\kappa$ media sobre las cinco variables), por nivel}
        \label{tab:iris-intraia}
    }{
        \begin{tabularx}{\textwidth}{
            >{\hsize=1.7\hsize\raggedright\arraybackslash}X
            >{\hsize=0.65\hsize\centering\arraybackslash}X
            >{\hsize=0.65\hsize\columncolor{gray!15}\centering\arraybackslash}X
        }
            \toprule
            \textbf{Par de modelos} & \boldmath$\kappa$ \textbf{(B0)} & \boldmath$\kappa$ \textbf{(B1)} \\
            \midrule
            \texttt{gemini} $\leftrightarrow$ \texttt{gpt-4o-mini}     & 0,079 & 0,052 \\
            \texttt{gemini} $\leftrightarrow$ \texttt{gpt-5.4-nano}    & 0,067 & 0,014 \\
            \texttt{gpt-4o-mini} $\leftrightarrow$ \texttt{gpt-5.4-nano} & \textbf{0,280} & 0,084 \\
            \midrule
            \texttt{gemma} $\leftrightarrow$ \texttt{gemini}          & 0,121 & \textbf{0,151} \\
            \texttt{gemma} $\leftrightarrow$ \texttt{gpt-4o-mini}     & 0,131 & 0,104 \\
            \texttt{gemma} $\leftrightarrow$ \texttt{gpt-5.4-nano}    & 0,111 & 0,025 \\
            \bottomrule
        \end{tabularx}
    }
    \endgroup
\end{table}

Los modelos concuerdan entre sí tan poco como con las personas ($\kappa$ entre
0,014 y 0,151 en B1). Un ``equipo de IA'' con criterio compartido no existe: cada
modelo tiene su propio patrón, y lo único que tienen en común es la dirección del
sesgo, no los casos concretos. Esto desmonta la lectura fácil de que ``la IA tiene
un sesgo'' uniforme. Hay cuatro sesgos distintos que solo coinciden en ser
permisivos. El modelo local \texttt{gemma4:e4b} es el que más solape de criterio
mantiene con el resto (sus pares llegan a 0,151 con \texttt{gemini}), pero ni
siquiera ese máximo alcanza un acuerdo apreciable. Repetir el análisis sin la
arquitectura de agentes (nivel B0) deja el patrón intacto, de modo que el sesgo es
del modelo y no del sistema de \textit{skills}. Un matiz sí cambia: las
\textit{Agent Skills} \textbf{reducen} el acuerdo entre modelos en cinco de los seis
pares (el par \texttt{gpt-4o-mini} y \texttt{gpt-5.4-nano} baja de 0,280 a 0,084),
lo que reitera que la arquitectura acentúa las diferencias de cada modelo en lugar
de acercarlos a un criterio común.


Los tres análisis convergen. La baja concordancia \textbf{no debe leerse como
incapacidad de los modelos para detectar sexismo}, sino como el producto de un
criterio experto amplio y en parte en disputa, de una heterogeneidad sustancial
entre anotadores (del mismo orden que la distancia entre humanas y máquinas, con
prevalencias medias que van de 0,283 a 0,639) y de la ambigüedad de algunas
definiciones del codebook. El coeficiente $\kappa$ mide, en buena medida, esa
distancia de criterio: \textbf{el techo lo pone la tarea, no el modelo}.




\section{Discusión}
\label{sec:discussion}
%----------
% DISCUSIÓN — lectura transversal de los resultados
%----------

Más allá de cada tabla, los resultados del capítulo admiten una lectura de
conjunto en torno a tres ideas: el carácter condicional de las \textit{Agent
Skills}, la lección metodológica que deja la evaluación pormenorizada que hemos hecho, y el techo que impone la
propia tarea, del que se derivan las implicaciones para el proyecto.

\subsection{El valor de las \textit{Agent Skills} es condicional al modelo}

El hallazgo central obliga a matizar una expectativa extendida: que dotar a un modelo
de más componentes (herramientas, conocimiento cargable bajo demanda, capacidad de
actuar) mejora por sí solo su desempeño. No es lo que se observa. Empaquetar la
metodología experta como \textit{skills} cargables no eleva el acuerdo de forma
automática, y lo que separa a los cuatro modelos no es tanto la dirección del cambio
como su magnitud. Solo \texttt{gemini} registra una ganancia de entidad, la de
\texttt{masc\_generico}, donde su $\kappa$ sube 0,188 puntos. En los demás los
movimientos se quedan en unas pocas milésimas o van hacia abajo, y
\texttt{gpt-5.4-nano} no mejora en ninguna de las cinco variables. La arquitectura no
es útil ni inútil en abstracto, sino que \textbf{amplifica lo que cada modelo ya
trae}. Aprovechar la carga bajo demanda exige que el modelo sepa decidir qué necesita
y cuándo, una competencia desigual entre modelos, y cuando falta, el componente aporta
contexto que satura en lugar de ayudar.

De la descomposición por componentes (Sección~\ref{subsec:iris_descomposicion}) se
sigue además una recomendación práctica, con la cautela de que ninguna configuración
se impone en las cinco variables. La recuperación en vivo sobre las guías no produce
una ganancia que justifique su coste, ya que sobre la configuración mínima mejora el
F1 de la clase positiva en nueve de las veinte celdas y lo empeora en ocho. La
preferible para un despliegue real es la de \textit{skills} con sus resúmenes, sin
herramienta, que encabeza siete de las veinte celdas, más que ninguna otra, y es a la
vez la más simple y la más barata. No es, con todo, una opción imbatible, ya que queda
por detrás de la configuración mínima en siete celdas, y en una de ellas por mucho, la
de \texttt{gpt-4o-mini} en \texttt{masc\_generico}, donde su F1 cae de 0,765 a
0,401.


\subsection{Exactitud, \texorpdfstring{$\kappa$}{kappa} y prevalencia: por qué informar de las tres}
\label{subsec:iris_tres_metricas}

El segundo aprendizaje es metodológico. En tareas con clases muy desequilibradas, la
exactitud y el coeficiente $\kappa$ pueden desacoplarse y llegar a moverse en
sentidos opuestos. Una intervención que vuelve al modelo menos conservador aumenta
los aciertos triviales por efecto de la distribución marginal, sin que mejore su capacidad
discriminativa. Este trabajo lo ilustra en dos celdas opuestas. En
\texttt{lenguaje\_sexista}, cuya clase mayoritaria es el \textit{Sí},
\texttt{gpt-4o-mini} se vuelve menos conservador al recibir las \textit{skills} y pasa
de marcar el 9,0\,\% de las piezas al 46,2\,\%, con lo que su exactitud sube de 0,236 a
0,457 mientras su $\kappa$ baja de $-0,013$ a $-0,037$ (Tabla~\ref{tab:var-v25}). En
\texttt{asimetria\_mujer\_hombre}, cuya clase mayoritaria es el \textit{No},
\texttt{gpt-5.4-nano} hace justo lo contrario y pasa del 37,5\,\% al 0,7\,\%, de modo
que su exactitud salta de 0,619 a 0,931 al tiempo que su recall cae de 0,457 a cero
(Tabla~\ref{tab:var-v33}). En ninguno de los dos casos discrimina mejor el modelo, sino
que se acerca a la respuesta que más se repite y la exactitud premia ese acercamiento.
Apoyarse solo en ella llevaría a concluir que las \textit{Agent Skills} mejoran el
sistema, cuando el $\kappa$ indica lo contrario. Por eso, en tareas
desequilibradas, la exactitud debe interpretarse siempre junto al $\kappa$ y a la
prevalencia, para no tomar un cambio en la distribución de respuestas por una mejora
real de criterio.


\subsection{Dos límites, uno de la tarea y otro de los modelos}
\label{subsec:iris_dos_limites}

El acuerdo absoluto es bajo, y el análisis del Capítulo permite repartir ese
desajuste entre dos causas que conviene no confundir.

La primera es un \textbf{límite de la tarea}. Dos equipos expertos codificaron la
misma variable con criterios muy distantes, hasta el punto de que
\texttt{sexismo\_discurso} pasa del 11,5\,\% al 65,2\,\% de piezas marcadas según
quién anote (Sección~\ref{subsec:iris_errores}). La brecha, además, no es
institucional. Dentro de un mismo equipo esa variable va del 28,3\,\% al 97,8\,\%
según quién codifique, un rango de 69,5 puntos que por sí solo supera a los 53,7 que
separan a los dos equipos, de modo que dos personas con la misma formación y el mismo
libro de códigos marcan, una, poco más de una de cada cuatro piezas, y la otra, casi
todas. Cuando el criterio admite esa
amplitud, el patrón de referencia deja de ser un suelo estable y el coeficiente
$\kappa$ mide en buena parte distancia entre criterios más que error del sistema.
Esta variabilidad no es un defecto del corpus de IRIS, sino un rasgo reconocido de
las tareas de anotación subjetiva, que la investigación reciente propone tratar
como información y no como ruido a eliminar~\cite{davani2022, plank2022}. La
campaña EXIST, referencia en identificación de sexismo en español, ha adoptado por
ello un paradigma de aprendizaje con desacuerdo en el que se conservan las
etiquetas de todas las personas anotadoras en lugar de agregarlas en una
sola~\cite{plaza2024}.

La segunda es un \textbf{límite de los modelos}, y sería un error disolverlo en el
anterior. En \texttt{lenguaje\_sexista} ambos equipos coinciden en marcar la
variable en la mayoría de las piezas, el 74,6\,\% y el 85,5\,\% respectivamente, y
aun así tres de los cuatro modelos presentan un recall prácticamente nulo. Ahí no
hay criterio en disputa, hay un fenómeno que las personas ven y el sistema no. El
caso examinado en la Sección~\ref{subsec:iris_caso} muestra el mecanismo, ya que
los modelos localizan el masculino genérico y lo declaran un colectivo mixto
válido, reproduciendo el mismo automatismo que se les pedía detectar.

El peso de cada límite cambia según la variable, y de ahí se sigue que el margen de
mejora tenga también dos vías. Una consiste en explicitar mejor el constructo en el
libro de códigos, apoyándose en anotaciones cualitativas que lo delimiten con
ejemplos, y conviene priorizarla en las variables de criterio más implícito. La
otra consiste en mejorar la sensibilidad del sistema en aquello que las expertas sí
identifican de forma concordante. Conviene además evitar una simplificación, ya que
los modelos no comparten un sesgo homogéneo, sino cuatro comportamientos distintos
que solo coinciden en su dirección.

\subsection{Implicaciones para IRIS}
\label{subsec:iris_implicaciones}

De los resultados se desprenden cuatro orientaciones prácticas para el proyecto.

\paragraph{El despliegue local es viable.} El modelo abierto \texttt{gemma},
ejecutado en infraestructura propia, compite con los servicios comerciales y los
supera en algunas variables, como \texttt{asimetria\_mujer\_hombre}, donde obtiene
el mejor recall y la mejor precisión de los cuatro
(Sección~\ref{subsec:iris_local}). Para una institución que maneja material sujeto
a privacidad, esto significa que preservar los datos y prescindir de tarifas no
obliga a renunciar al rendimiento, a cambio de asumir un coste de cómputo. Ese coste,
además, se comporta al revés que el de los modelos de pago, en los que el empaquetado
multiplica la factura por 2,2. En el local la reduce, ya que procesar una pieza pasa
de 319,3 segundos en el nivel de control a 218,6 en el desplegado, porque la carga bajo
demanda incorpora solo la \textit{skill} que hace falta en lugar de arrastrar la
metodología completa en cada llamada.

\paragraph{Conviene tratar cada variable por separado.} Ni la configuración óptima
ni la utilidad del sistema son las mismas en las cinco. En
\texttt{masc\_generico} el sistema detecta ya casi tres de cada cuatro casos
señalados por las expertas y puede incorporarse al flujo de trabajo, mientras que en
\texttt{lenguaje\_sexista} no aporta información utilizable pese a ser la variable
que más presupuesto de inferencia consume, el 24,2\,\% del total
(Sección~\ref{subsec:iris_coste}).
Concentrar el esfuerzo donde el sistema funciona, y no distribuirlo por igual,
mejora el rendimiento y reduce el gasto a la vez. Esa afirmación admite una cifra.
Asignar a cada variable el modelo que mejor se comporta en ella supone 8,77 USD sobre
las 1.313 piezas, cantidad en la que \texttt{asimetria\_mujer\_hombre} no devenga
tarifa por recaer en el modelo local, si bien consume tiempo de cómputo. Fijar en
cambio un único modelo para las cinco costaría 13,16 USD con \texttt{gemini}, que es el
de mayor acuerdo, y 6,89 USD con \texttt{gpt-4o-mini}, que sale más barato pero queda
por detrás en tres de las cinco variables. La asignación por variable se sitúa, por
tanto, entre ambas en coste y por delante de las dos en rendimiento, que es cuanto cabe
pedirle a una decisión de despliegue.

\paragraph{La supervisión humana no es un remedio provisional.} Que el sistema
requiera revisión experta no constituye una anomalía de esta implementación, sino
lo que la literatura viene recomendando para las tareas de anotación con carga
interpretativa. Pangakis et al.~\cite{pangakis2023} replican veintisiete tareas de
anotación procedentes de la investigación en ciencias sociales y concluyen que
ningún proceso automático puede darse por válido sin contrastarlo con etiquetas
humanas, ya que en nueve de ellas la precisión o el recall caen por debajo de 0,5.
Kocoń et al.~\cite{kocon2023} constatan que el rendimiento de un modelo general se
degrada de forma acusada justo en las tareas subjetivas, con una pérdida media del
25,5\,\% frente a los sistemas especializados. Y desde la perspectiva de la
gobernanza, Gorwa et al.~\cite{gorwa2020} advierten de que los sistemas
algorítmicos de moderación, aun bien optimizados, siguen siendo opacos y difíciles
de fiscalizar. El paradigma \textit{human-in-the-loop} que adopta este trabajo se
sitúa, por tanto, en el consenso de la disciplina.

\paragraph{Ante la duda, conviene marcar de más.} Como herramienta de apoyo a la
postedición, el sistema no decide, sino que selecciona lo que merece una segunda
lectura. En ese uso los dos errores no cuestan lo mismo, ya que un falso positivo
supone unos segundos de descarte para la analista mientras que un falso negativo
es un caso que nadie revisará (Sección~\ref{sec:metricas_eval}). Los resultados
apuntan en la dirección contraria a la deseable, puesto que el sistema falla
precisamente por defecto, de modo que cualquier ajuste futuro debería orientarse a
elevar el recall aun a costa de la precisión. La Sección~\ref{subsec:iris_equipo_ia}
apunta una vía concreta en ese sentido, la de marcar la pieza cuando lo haga
cualquiera de los modelos, que eleva el recall en las cinco variables, si bien obliga
a ejecutarlos todos sobre cada pieza.

Con estas cuatro orientaciones, el sistema encaja mejor como asistente que como
decisor. Su salida trazable, con código, explicación y evidencias literales,
permite que la analista audite, acepte o corrija cada predicción, y el caso
analizado en la Sección~\ref{subsec:iris_caso} muestra que esa explicación no es un
añadido cosmético, sino lo que permite entender por qué falla el sistema y, con
ello, corregirlo. El valor no está en sustituir el criterio experto, sino en
ordenar el corpus, señalar los casos que merecen atención y hacerlo de forma
auditable.



%----------
%	CONCLUSIONES Y LÍNEAS FUTURAS
%----------

\chapter{Conclusiones y Líneas Futuras}
%----------
%	CONCLUSIONS AND FUTURE WORK
%----------	

This Chapter summarizes the main conclusions of this work, based on the analysis and results discussed in Chapter~\ref{results_and_discussion}. Moreover, we would also take another look and expand the applications mentioned in Section~\ref{sec:socioeconomic_impact}, exploring ideas for future work and potential use cases, and providing insights on how this research could be further developed and used in real-world scenarios.


\section{Conclusiones}
\label{sec:conclusions}
%----------
%	CONCLUSIONES
%----------

Este trabajo ha abordado la detección automática de lenguaje sexista en la prensa
española escrita, a partir del corpus y la metodología del proyecto IRIS, con el fin de asistir el análisis
experto de la representación mediática de las mujeres. La propuesta central ha sido
empaquetar la metodología experta como \textit{Agent Skills} cargables bajo demanda,
sobre una arquitectura de agentes especializadas, una por cada variable del libro de
códigos (V25, V26, V30, V33 y V35). Para valorar su aportación real, esa arquitectura
(nivel B1) se comparó con una condición de control que recibía la misma metodología
inyectada en el \textit{prompt} (nivel B0), sobre un corpus de 1.313 noticias anotadas
y cuatro modelos de lenguaje, tres accesibles por API y uno local.

La evaluación experimental, interpretada en detalle en la Sección~\ref{sec:discussion},
deja varias conclusiones que aquí se recogen de forma sintética. La primera es que
empaquetar la metodología como \textit{Agent Skills} \textbf{no mejora el acuerdo por
sí mismo}, ya que solo un modelo registra una ganancia de entidad, la de
\texttt{gemini} en \texttt{masc\_generico}, donde su $\kappa$ sube 0,188 puntos,
mientras en los demás los movimientos se quedan en unas pocas milésimas o van hacia
abajo. La arquitectura amplifica, por tanto, la capacidad que cada modelo ya trae en
lugar de conferirla, y la herramienta de recuperación en vivo (RAG) no aporta una
ganancia que justifique su coste. La segunda es de alcance metodológico: en una tarea con
clases muy desequilibradas, la exactitud y $\kappa$ pueden moverse en sentidos opuestos,
por lo que reportar solo la exactitud llevaría a la conclusión contraria y se impone
informar de exactitud, $\kappa$ y prevalencia de forma conjunta. La tercera es que el
bajo acuerdo absoluto responde a \textbf{dos causas que conviene no confundir}. Una es
un techo de la propia tarea, documentado en \texttt{sexismo\_discurso}, que dos equipos
marcan en el 11,5\,\% y en el 65,2\,\% de las piezas y que dentro de un mismo equipo va
del 28,3\,\% al 97,8\,\%, reflejo de lo difícil que resulta percibir el fenómeno incluso
a quien lo busca. La otra es una limitación de los modelos: en \texttt{lenguaje\_sexista} los dos
equipos marcan la variable en la mayoría de sus piezas, el 74,6\,\% y el 85,5\,\%
respectivamente, y aun así tres de los cuatro modelos apenas la detectan, con un
recall prácticamente nulo. El peso de cada una cambia según la variable, y reparte responsabilidades, de
manera que el desacuerdo recoge a la vez la dificultad de percibir el fenómeno y el
error de detección del sistema. Cabe
precisar, además, que los modelos quedan, salvo \texttt{gemini} en
\texttt{masc\_generico}, por debajo de la persona que menos marca en términos de \textit{recall}, lo que responde a
una capacidad de detección menor y no a un criterio distinto. La cuarta es que un modelo local
de pesos abiertos, \texttt{gemma4:e4b}, compite variable a variable con los servicios
comerciales accesibles por API, a los que supera en \texttt{asimetria\_mujer\_hombre}, con la cautela de que esa
ventaja descansa en un solo acierto, y a dos de los cuales adelanta en
\texttt{masc\_generico}, preservando la privacidad del material y sin
tarifas de inferencia, a cambio de trasladar el coste al tiempo de cómputo. Realizar esa comparación sobre la misma tarea y el mismo corpus, con un flujo de procesamiento agnóstico a la
empresa proveedora, ha sido uno de los ejes metodológicos del trabajo.

En cuanto al cumplimiento de los objetivos de la Sección~\ref{sec:objectives_motivation},
el balance es el siguiente. El objetivo general (OG1) se ha satisfecho al investigar y
evaluar de forma sistemática una arquitectura basada en agentes y LLM para detectar
variables cuantitativas y cualitativas de representación mediática, delimitando con evidencia empírica
tanto sus prestaciones como sus límites. El diseño conceptual del sistema, con agentes
especializadas y sus flujos de análisis y verificación (OE1), y las estrategias de
\textit{prompting} y herramientas asociadas, criterios, plantillas y contextos expertos
empaquetados como \textit{skills} (OE2), se desarrollaron en el
Capítulo~\ref{implementation} y se validaron en el Capítulo~\ref{results_and_discussion}.
Las técnicas de representación y recuperación de información, incluida la RAG y la
similitud léxica (OE3), se implementaron y evaluaron, con el hallazgo de que la recuperación
en vivo no produce una ganancia consistente que justifique su coste. La comparación
entre modelos y su idoneidad para detectar sesgos explícitos e implícitos (OE4)
constituye el núcleo del Capítulo~\ref{results_and_discussion}. Por último, la propuesta
de directrices para la integración editorial y la supervisión humana (OE5) se concreta
en las directrices que se enuncian al cierre de este capítulo y en las líneas futuras
que se detallan después.

De todo ello se desprende, en coherencia con la discusión, que el sistema encaja mejor
como \textbf{asistente que como decisor}. Su salida tiene dos planos, y solo uno se ha
podido evaluar. El primero es la etiqueta, el plano cuantitativo, que se contrasta con
la anotación de referencia y es el que se mide aquí. El segundo es el marcado de
evidencias y la explicación del veredicto, es decir, el plano cualitativo, esa segunda dimensión
que un análisis de contenido experto también necesita. Este último no cuenta todavía
con un marcado experto de referencia con el que compararse, de modo que su validación
queda como línea futura. El valor del sistema es, en definitiva, práctico y subordinado
al criterio experto: ordenar el corpus, señalar los casos que merecen una segunda
lectura y aliviar el trabajo repetitivo de quienes codifican.

El alcance de estas conclusiones está acotado por varias limitaciones. El corpus
procede de la prensa española y se analizan cinco de las quince variables de lenguaje
del libro de códigos, de modo que la generalización a otros medios, idiomas o variables
queda por comprobar. La evaluación se ha realizado en régimen \textit{zero-shot} y
sobre modelos ligeros y económicos, sin explorar el ajuste fino (\textit{fine-tuning})
ni modelos de mayor capacidad. La recuperación sobre las guías se implementó con similitud léxica (TF-IDF),
y una variante semántica podría comportarse de otro modo. Dos de las variables
descansan además sobre muy pocos casos positivos, 81 en
\texttt{asimetria\_mujer\_hombre} y 131 en \texttt{denominacion\_sexualizada}, de modo
que sus cifras se juegan en un puñado de piezas y admiten menos confianza que las
demás. Tampoco ha podido evaluarse el salto semántico, ya que la anotación de
referencia no emplea ese valor en ninguna pieza del corpus. Como cada pieza fue
codificada por una sola persona, no fue posible calcular la fiabilidad
intercodificadora, por lo que el techo de la tarea se acota de forma indirecta y con incertidumbre significativa.
Por último, y como ya se ha señalado, solo se ha evaluado el plano cuantitativo de la
salida. Estas limitaciones no invalidan los hallazgos, pero delimitan su alcance y
orientan las líneas de trabajo que se proponen a continuación.

De cara a una futura integración del sistema en entornos editoriales reales (OE5), el
trabajo permite formular siete directrices:
\begin{itemize}
    \item Desplegarlo como asistente que ordena y prioriza el corpus, y no como decisor
    automático.
    \item Preservar la trazabilidad de cada predicción, con sus evidencias y explicación,
    para que la analista pueda auditarla, aceptarla o corregirla.
    \item Mantener a la persona en el bucle, de modo que su criterio prevalezca sobre la
    salida del modelo, en especial en las variables en las que el sistema deja sin
    señalar la mayor parte de los positivos.
    \item Decidir su uso variable a variable, asignando a cada una el modelo que mejor se
    comporta en ella, en lugar de desplegar una única configuración para las cinco.
    \item Orientar cualquier ajuste futuro a elevar el \textit{recall}, aun a costa de la
    precisión, porque en postedición un falso negativo es un caso que nadie revisará y
    un falso positivo cuesta unos segundos de descarte.
    \item Favorecer el despliegue local cuando el material sea sensible, por razones de
    privacidad y de coste.
    \item Informar del desempeño con métricas conscientes del desequilibrio, esto es, acuerdo
    y prevalencia además de exactitud, para no sobrestimar la fiabilidad del sistema.
\end{itemize}


\section{Líneas Futuras}
\label{sec:futurework_usecases}
%----------
%	LÍNEAS FUTURAS
%----------

Los resultados de este trabajo abren varias vías de continuación, ordenadas de la más
inmediata a la más ambiciosa.

La primera se deriva directamente del principal hallazgo. Si buena parte del desacuerdo
mide distancia de criterio y no error, tiene poco sentido seguir tratando cada variable
como si tuviera una única interpretación correcta. Resultaría prometedor explorar
enfoques que no asuman una etiqueta de oro única por categoría, como el aprendizaje con
desacuerdo y las aproximaciones multietiqueta, que permitirían al sistema representar la
ambigüedad de las variables más disputadas en lugar de forzar un veredicto binario. En
la misma línea, incorporar una doble codificación de una parte del corpus, ausente en
este estudio (Sección~\ref{subsec:limitacion_fiabilidad}), permitiría medir la
fiabilidad intercodificadora humana y situar con precisión el techo real de la tarea,
hoy solo acotado de forma indirecta.

Una segunda vía es refinar el aprovechamiento del contexto experto por parte de los
modelos. Dado que la ganancia de las \textit{Agent Skills} depende de la capacidad del
modelo para explotar la carga bajo demanda, conviene estudiar estrategias de
\textit{few-shot} con ejemplos verificados, así como el ajuste fino (\textit{fine-tuning})
de modelos abiertos como \texttt{gemma4:e4b} sobre el propio corpus anotado. El modelo
local, competitivo en acuerdo y sin coste de tarifas
(Sección~\ref{subsec:iris_local}), es el candidato natural para esta especialización,
con la ventaja añadida de preservar la privacidad del material.

Una tercera vía atañe a la explicabilidad y a la validación del plano cualitativo. El
sistema ya produce, para cada predicción, la explicación y las evidencias literales que
la sustentan, pero ese marcado no ha podido evaluarse por falta de un marcado experto de
referencia. Construir ese conjunto de referencia, en el que las analistas señalen sobre
el texto los fragmentos que motivan cada variable, permitiría medir la calidad de las
evidencias, no solo de la etiqueta, y cerrar así la evaluación de las dos dimensiones que
un análisis de contenido experto necesita. Sobre esa misma base trazable cabe construir
un flujo con humano en el bucle, en el que la codificadora acepte, rechace o corrija cada
veredicto y esa retroalimentación sirva para reajustar los criterios y los umbrales por
variable, de forma coherente con el papel de asistente, y no de decisor, que se ha
argumentado para el sistema.

Por último, en el plano de la aplicación, el objetivo a medio plazo es la integración
del sistema en una herramienta de postedición dentro del proyecto IRIS, que ordene el
corpus, resalte los pasajes candidatos y facilite el trabajo de las analistas. Este
trabajo se enmarca en la colaboración del proyecto IRIS con RTVE, cuyo Observatorio de
Igualdad aporta parte de las guías de lenguaje que fundamentan el libro de códigos, de
modo que un escenario de aplicación natural es el apoyo a los flujos de revisión
editorial de una corporación de radiotelevisión pública, donde el seguimiento de la
representación de las mujeres forma parte de su mandato de servicio público. Esa
herramienta permitiría además extender el análisis a más variables del libro de códigos
y a otros medios y periodos, y ofrecería a instituciones y equipos de investigación un
instrumento reproducible para el seguimiento sistemático de la representación de las
mujeres en los medios. La combinación de un despliegue local respetuoso con la privacidad
y una salida auditable por personas expertas es, a nuestro juicio, la vía más realista
para llevar este trabajo del entorno experimental al uso real.


%----------
%	BIBLIOGRAFÍA
%----------	

%\nocite{*} % Si quieres que aparezcan en la bibliografía todos los documentos que la componen (también los que no estén citados en el texto) descomenta está lína

\clearpage
\addcontentsline{toc}{chapter}{Bibliografía}

\printbibliography



%----------
%	ANEXOS
%----------	
\appendix
\chapter{Presupuesto}
\label{appendix_budget}
\pagenumbering{gobble} % Las páginas de los anexos no se numeran
%----------
%	APPENDIX 1: BUDGET
%----------	

En este anexo se desglosa la estimación económica necesaria para el desarrollo del presente TFM, cuyos importes quedan recogidos en la Tabla~\ref{tab:budget}.


En cuanto a los costes operativos de personal, se ha cuantificado la dedicación del estudiante a razón de 5 horas diarias a lo largo de los 6 meses que ha durado el proyecto. Tomando como referencia el contrato de investigación de IRIS\_IAMEDIA, supone una tarifa total aproximada de 9.300 € a razón de 15,5 €/hora. Por su parte, la supervisión y dirección académica ha corrido a cargo de dos tutoras, con las que se han mantenido en torno a 25 reuniones de una hora a lo largo de los seis meses del proyecto, esto es, aproximadamente una reunión semanal. Al asistir ambas tutoras, la dedicación conjunta asciende a unas 50 horas. Asignando a este rol de dirección una tarifa de 25,64 € por hora (equivalente a un perfil de ingeniería senior o investigadora principal\footnote{Fuente: \url{https://es.talent.com/salary?job=senior+data+analyst}, [Último acceso: 28/07/2026]}), esta partida suma un total de 1.282,00 €.


En cuanto al equipamiento informático, el desarrollo se ha llevado a cabo utilizando un equipo MacBook Pro M4 de 2024, cuyo valor de mercado se sitúa en 2.500 €. Siguiendo los criterios contables estandarizados, la vida útil de este dispositivo se amortiza en un periodo de 4 años, lo que se traduce en una depreciación de 625 € anuales. Al circunscribirse el uso de la máquina a los 6 meses de duración de esta investigación, la cantidad imputable al presupuesto general es de 312,50 €. Adicionalmente, la ejecución del modelo local \texttt{gemma4:e4b} se realizó sobre la granja de GPU (NVIDIA RTX 3090 y 4090) del Departamento de Teoría de la Señal y Comunicaciones (TSC) de la Universidad Carlos III de Madrid. Al tratarse de infraestructura institucional de investigación, su uso no ha supuesto un coste directo para el proyecto y se imputa a 0 €, si bien su disponibilidad ha sido esencial para el despliegue agéntico del modelo local.

Finalmente, el apartado de software contempla tres partidas: 90 € para el acceso a las noticias de pago de los periódicos analizados, 131,94 € para la suscripción al modelo Gemini Pro (Gemini Advanced) durante los seis meses, y el coste de inferencia de las API comerciales empleadas en la evaluación, esto es, Gemini de Google y los modelos de OpenAI, que ascendió a unos 136 USD (en torno a 126,60 €). Esa cifra incluye las ejecuciones de los dos niveles comparados en el Capítulo~\ref{results_and_discussion} y las tres configuraciones adicionales de la ablación de la Sección~\ref{subsec:iris_descomposicion}, mientras que las cifras de coste del Capítulo~\ref{results_and_discussion} recogen únicamente los dos niveles. Conviene matizar que el acceso a los periódicos y la suscripción a Gemini Pro no han supuesto un desembolso real, al estar cubiertos por las licencias y suscripciones institucionales de la Universidad Carlos III de Madrid. El coste de inferencia de las API, en cambio, sí ha constituido un gasto efectivo del proyecto. De este modo, la ejecución global del proyecto supone un coste estimado de 11.243,04 €.




\begin{table}[H]
    \centering
    \caption{Presupuesto para la investigación}
    \label{tab:budget}
    \small % font size
    \begin{tabularx}{\textwidth}{
        >{\hsize=0.75\hsize}X % Category
        >{\hsize=0.75\hsize}X % Name
        >{\hsize=2.25\hsize}X % Details
        >{\hsize=0.625\hsize\centering\arraybackslash}X % Period
        >{\hsize=0.625\hsize\centering\arraybackslash}X % Cost
        }
        \hline
        \rowcolor{gray!30}
        \textbf{Categoría} & \textbf{Nombre} & \textbf{Detalles} & \textbf{Periodo} & \textbf{Coste (€)} \\
        \hline
        \multirow{1}{*}{Coste Humano} & Estudiante & 5 horas/día (15,5 €/h) & 6 meses & 9.300,00 \\
        \cline{2-5}
        & Tutoras (2) & 25 reuniones $\times$ 2 tutoras = 50 h (25,64 €/h) & 6 meses & 1.282,00 \\
        %\cline{2-5}
        %& Expert & Expert consultation and annotation, 4 meetings (25.64€/hour) & 1 month & 769.2 \\
        \hline
        \multirow{2}{*}{Hardware} & MacBook Pro M4 (2024) & 2.500 € de valor, amortizado en 4 años (625 €/año) & 6 meses & 312,50 \\
        \cline{2-5}
        & Granja de GPU (TSC) & NVIDIA RTX 3090/4090, infraestructura institucional & 6 meses & 0,00 \\
        \hline
        \multirow{3}{*}{Software} & Licencias periódicos & Suscripción digital básica & 6 meses & 90,00 \\
        \cline{2-5}
        & Gemini Pro & Suscripción a Gemini Advanced & 6 meses & 131,94 \\
        \cline{2-5}
        & API de inferencia & Gemini (Google) y OpenAI, $\approx$136 USD & 6 meses & 126,60 \\
        \hline
        \multicolumn{4}{r}{\textbf{Total}} & \textbf{11.243,04} \\
        \hline
    \end{tabularx}
\end{table}



\chapter{Código}
\label{appendix_code}
\pagenumbering{gobble} % Las páginas de los anexos no se numeran
%----------
%	APÉNDICE 2: CÓDIGO FUENTE
%----------

El código desarrollado para este Trabajo de Fin de Máster comprende la
implementación de la arquitectura de agentes especializadas y de las
\textit{Agent Skills} descritas en el Capítulo~\ref{implementation}, junto con los
\textit{scripts} de ejecución y evaluación empleados en el
Capítulo~\ref{results_and_discussion}.

En lugar de reproducir el código en esta memoria, se ha optado por depositarlo en
un repositorio privado\footnote{\url{https://github.com/jgarcela/tfm-iris-uc3m}
[Último acceso: 14/09/2026].}, de manera que pueda consultarse en su forma
ejecutable. El acceso se facilita a petición al tribunal y a las personas del
proyecto que lo soliciten.

El repositorio contiene únicamente el código. El corpus periodístico y las
anotaciones expertas pertenecen a los proyectos IRIS\_IAMEDIA e InfoIA y no se
difunden con él, de acuerdo con las condiciones de uso descritas en la
Sección~\ref{sec:reg_framework}. El repositorio se mantiene privado por esa
misma razón, ya que las \textit{skills} incorporan el libro de códigos del proyecto.


\chapter{AI Act Compliance}
\label{appendix_ai_act}
\pagenumbering{gobble} % Las páginas de los anexos no se numeran
%----------
%	APPENDIX 3: AI ACT
%----------	

\begin{figure}[H]
	\ffigbox[\FBwidth]
	{\caption{AI Act Compliance Checker Responses for Our Project}\label{fig:appendix_ai_act}}
 % [Aquí ponemos como queremos que aparezca en el índice, sin la referencia]{Aquí como queremos que aparezca en el documento, con la referencia}
	{\includegraphics[scale=0.75]{Plantilla_TFG_ingles_2019/imagenes/ai_act_results.png}}
\end{figure}

\begin{figure}[H]
	\ffigbox[\FBwidth]
	{\caption{AI Act Compliance Checker Results for Our Project}\label{fig:appendix_ai_act_2}}
 % [Aquí ponemos como queremos que aparezca en el índice, sin la referencia]{Aquí como queremos que aparezca en el documento, con la referencia}
	{\includegraphics[scale=0.45]{Plantilla_TFG_ingles_2019/imagenes/ai_act_results_2.png}}
\end{figure}


\end{document}